
\Data\ID{1003108701}
\Updated{2020-07-15 03:07:21}
\Level{A}
\SubArea{Infinite series}
\LANGen{
\Question{
Use summation notation to express the given infinite series.
\[1+\frac12+\frac14+\frac18+\dots \]}
{\( \sum\limits_{n=1}^{\infty} \frac1{2^{n-1}} \)}{\( \sum\limits_{n=1}^{\infty} \frac1{2^{n}} \)}{\( \sum\limits_{n=1}^{\infty} \frac1{2^{n+1}} \)}{\( \sum\limits_{n=1}^{\infty} \frac1{2^{2n}} \)}{\( \sum\limits_{n=1}^{\infty} \frac1{2^{2n-1}} \)}{}}
\LANGpl{
\Question{
Użyj uproszczonego zapisu, aby wyrazić podane szeregi nieskończone. 
\[1+\frac12+\frac14+\frac18+\dots \]}
{\( \sum\limits_{n=1}^{\infty} \frac1{2^{n-1}} \)}{\( \sum\limits_{n=1}^{\infty} \frac1{2^{n}} \)}{\( \sum\limits_{n=1}^{\infty} \frac1{2^{n+1}} \)}{\( \sum\limits_{n=1}^{\infty} \frac1{2^{2n}} \)}{\( \sum\limits_{n=1}^{\infty} \frac1{2^{2n-1}} \)}{}}
\LANGcs{
\Question{
Danou nekonečnou geometrickou řadu zapište pomocí sumy:
\[1+\frac12+\frac14+\frac18+\dots \]}
{\( \sum\limits_{n=1}^{\infty} \frac1{2^{n-1}} \)}{\( \sum\limits_{n=1}^{\infty} \frac1{2^{n}} \)}{\( \sum\limits_{n=1}^{\infty} \frac1{2^{n+1}} \)}{\( \sum\limits_{n=1}^{\infty} \frac1{2^{2n}} \)}{\( \sum\limits_{n=1}^{\infty} \frac1{2^{2n-1}} \)}{}}
\LANGsk{
\Question{
Daný nekonečný geometrický rad zapíšte pomocou sumy:
\[1+\frac12+\frac14+\frac18+\dots \]}
{\( \sum\limits_{n=1}^{\infty} \frac1{2^{n-1}} \)}{\( \sum\limits_{n=1}^{\infty} \frac1{2^{n}} \)}{\( \sum\limits_{n=1}^{\infty} \frac1{2^{n+1}} \)}{\( \sum\limits_{n=1}^{\infty} \frac1{2^{2n}} \)}{\( \sum\limits_{n=1}^{\infty} \frac1{2^{2n-1}} \)}{}}
\LANGes{
\Question{
Usa la notación sumatoria para representar la serie infinita dada:
\[1+\frac12+\frac14+\frac18+\dots \]}
{\( \sum\limits_{n=1}^{\infty} \frac1{2^{n-1}} \)}{\( \sum\limits_{n=1}^{\infty} \frac1{2^{n}} \)}{\( \sum\limits_{n=1}^{\infty} \frac1{2^{n+1}} \)}{\( \sum\limits_{n=1}^{\infty} \frac1{2^{2n}} \)}{\( \sum\limits_{n=1}^{\infty} \frac1{2^{2n-1}} \)}{}}
\End

\Data\ID{1003108702}
\Updated{2020-07-15 03:07:19}
\Level{A}
\SubArea{Infinite series}
\LANGen{
\Question{
Use summation notation to express the given infinite series.
\[ -1+2-4+8-16+\dots \]}
{\( \sum\limits_{n=1}^{\infty} (-1)^n\cdot2^{n-1} \)}{\( \sum\limits_{n=1}^{\infty} (-1)^{n-1}\cdot2^{n-1} \)}{\( \sum\limits_{n=1}^{\infty} (-1)^{n+1}\cdot2^{n-1} \)}{\( \sum\limits_{n=1}^{\infty} (-1)^n\cdot2^{n+1} \)}{\( \sum\limits_{n=1}^{\infty} (-1)^n\cdot2^{n} \)}{}}
\LANGpl{
\Question{
Użyj uproszczonego zapisu, aby wyrazić podane szeregi nieskończone. 
\[ -1+2-4+8-16+\dots \]}
{\( \sum\limits_{n=1}^{\infty} (-1)^n\cdot2^{n-1} \)}{\( \sum\limits_{n=1}^{\infty} (-1)^{n-1}\cdot2^{n-1} \)}{\( \sum\limits_{n=1}^{\infty} (-1)^{n+1}\cdot2^{n-1} \)}{\( \sum\limits_{n=1}^{\infty} (-1)^n\cdot2^{n+1} \)}{\( \sum\limits_{n=1}^{\infty} (-1)^n\cdot2^{n} \)}{}}
\LANGcs{
\Question{
Danou nekonečnou geometrickou řadu zapište pomocí sumy:
\[ -1+2-4+8-16+\dots \]}
{\( \sum\limits_{n=1}^{\infty} (-1)^n\cdot2^{n-1} \)}{\( \sum\limits_{n=1}^{\infty} (-1)^{n-1}\cdot2^{n-1} \)}{\( \sum\limits_{n=1}^{\infty} (-1)^{n+1}\cdot2^{n-1} \)}{\( \sum\limits_{n=1}^{\infty} (-1)^n\cdot2^{n+1} \)}{\( \sum\limits_{n=1}^{\infty} (-1)^n\cdot2^{n} \)}{}}
\LANGsk{
\Question{
Daný nekonečný geometrický rad zapíšte pomocou sumy:
\[ -1+2-4+8-16+\dots \]}
{\( \sum\limits_{n=1}^{\infty} (-1)^n\cdot2^{n-1} \)}{\( \sum\limits_{n=1}^{\infty} (-1)^{n-1}\cdot2^{n-1} \)}{\( \sum\limits_{n=1}^{\infty} (-1)^{n+1}\cdot2^{n-1} \)}{\( \sum\limits_{n=1}^{\infty} (-1)^n\cdot2^{n+1} \)}{\( \sum\limits_{n=1}^{\infty} (-1)^n\cdot2^{n} \)}{}}
\LANGes{
\Question{
Usa la notación sumatoria para representar la serie infinita dada:
\[ -1+2-4+8-16+\dots \]}
{\( \sum\limits_{n=1}^{\infty} (-1)^n\cdot2^{n-1} \)}{\( \sum\limits_{n=1}^{\infty} (-1)^{n-1}\cdot2^{n-1} \)}{\( \sum\limits_{n=1}^{\infty} (-1)^{n+1}\cdot2^{n-1} \)}{\( \sum\limits_{n=1}^{\infty} (-1)^n\cdot2^{n+1} \)}{\( \sum\limits_{n=1}^{\infty} (-1)^n\cdot2^{n} \)}{}}
\End

\Data\ID{1003108703}
\Updated{2020-07-15 03:07:19}
\Level{A}
\SubArea{Infinite series}
\LANGen{
\Question{
Use summation notation to express the given infinite series.
\[ \frac3{x^3}+\frac3{x^2}+\frac3x+3+3x+\dots \]}
{\( \sum\limits_{n=1}^{\infty}3\cdot x^{n-4} \)}{\( \sum\limits_{n=1}^{\infty}3\cdot x^{n-3} \)}{\( \sum\limits_{n=1}^{\infty}3\cdot x^{n+3} \)}{\( \sum\limits_{n=1}^{\infty}3\cdot x^{n+4} \)}{\( \sum\limits_{n=1}^{\infty}3\cdot x^{n} \)}{}}
\LANGpl{
\Question{
Użyj uproszczonego zapisu, tak by wyrazić podane szeregi nieskończone. 
\[ \frac3{x^3}+\frac3{x^2}+\frac3x+3+3x+\dots \]}
{\( \sum\limits_{n=1}^{\infty}3\cdot x^{n-4} \)}{\( \sum\limits_{n=1}^{\infty}3\cdot x^{n-3} \)}{\( \sum\limits_{n=1}^{\infty}3\cdot x^{n+3} \)}{\( \sum\limits_{n=1}^{\infty}3\cdot x^{n+4} \)}{\( \sum\limits_{n=1}^{\infty}3\cdot x^{n} \)}{}}
\LANGcs{
\Question{
Danou nekonečnou geometrickou řadu zapište pomocí sumy:
\[ \frac3{x^3}+\frac3{x^2}+\frac3x+3+3x+\dots \]}
{\( \sum\limits_{n=1}^{\infty}3\cdot x^{n-4} \)}{\( \sum\limits_{n=1}^{\infty}3\cdot x^{n-3} \)}{\( \sum\limits_{n=1}^{\infty}3\cdot x^{n+3} \)}{\( \sum\limits_{n=1}^{\infty}3\cdot x^{n+4} \)}{\( \sum\limits_{n=1}^{\infty}3\cdot x^{n} \)}{}}
\LANGsk{
\Question{
Daný nekonečný geometrický rad zapíšte pomocou sumy:
\[ \frac3{x^3}+\frac3{x^2}+\frac3x+3+3x+\dots \]}
{\( \sum\limits_{n=1}^{\infty}3\cdot x^{n-4} \)}{\( \sum\limits_{n=1}^{\infty}3\cdot x^{n-3} \)}{\( \sum\limits_{n=1}^{\infty}3\cdot x^{n+3} \)}{\( \sum\limits_{n=1}^{\infty}3\cdot x^{n+4} \)}{\( \sum\limits_{n=1}^{\infty}3\cdot x^{n} \)}{}}
\LANGes{
\Question{
Usa la notación sumatoria para representar la serie infinita dada:
\[ \frac3{x^3}+\frac3{x^2}+\frac3x+3+3x+\dots \]}
{\( \sum\limits_{n=1}^{\infty}3\cdot x^{n-4} \)}{\( \sum\limits_{n=1}^{\infty}3\cdot x^{n-3} \)}{\( \sum\limits_{n=1}^{\infty}3\cdot x^{n+3} \)}{\( \sum\limits_{n=1}^{\infty}3\cdot x^{n+4} \)}{\( \sum\limits_{n=1}^{\infty}3\cdot x^{n} \)}{}}
\End

\Data\ID{1003108704}
\Updated{2021-05-05 06:05:28}
\Level{A}
\SubArea{Infinite series}
\LANGen{
\Question{
Expand the given sum by listing individual terms.
\[ \sum\limits_{n=2}^6(4n-5) \]}
{\( 3+7+11+15+19 \)}{\( -1+3+7+11+15 \)}{\( 4 + 8 +12 +16 +20 \)}{\( 3+ 8 + 11 + 15 + 19 \)}{\( 8 - 5 + 12 + 16 + 20 \)}{}}
\LANGpl{
\Question{
Rozszerz podane sumy wypisując wszystkie wyrazy tego szeregu. 
\[ \sum\limits_{n=2}^6(4n-5) \]}
{\( 3+7+11+15+19 \)}{\( -1+3+7+11+15 \)}{\( 4 + 8 +12 +16 +20 \)}{\( 3+ 8 + 11 + 15 + 19 \)}{\( 8 - 5 + 12 + 16 + 20 \)}{}}
\LANGcs{
\Question{
Danou sumu rozepište pomocí součtu:
\[ \sum\limits_{n=2}^6(4n-5) \]}
{\( 3+7+11+15+19 \)}{\( -1+3+7+11+15 \)}{\( 4 + 8 +12 +16 +20 \)}{\( 3+ 8 + 11 + 15 + 19 \)}{\( 8 - 5 + 12 + 16 + 20 \)}{}}
\LANGsk{
\Question{
Danú sumu rozpíšte pomocou súčtu:
\[ \sum\limits_{n=2}^6(4n-5) \]}
{\( 3+7+11+15+19 \)}{\( -1+3+7+11+15 \)}{\( 4 + 8 +12 +16 +20 \)}{\( 3+ 8 + 11 + 15 + 19 \)}{\( 8 - 5 + 12 + 16 + 20 \)}{}}
\LANGes{
\Question{
Elige el sumatorio adecuado para la serie dada :
\[ \sum\limits_{n=2}^6(4n-5) \]}
{\( 3+7+11+15+19 \)}{\( -1+3+7+11+15 \)}{\( 4 + 8 +12 +16 +20 \)}{\( 3+ 8 + 11 + 15 + 19 \)}{\( 8 - 5 + 12 + 16 + 20 \)}{}}
\End

\Data\ID{1003108705}
\Updated{2021-05-05 06:05:42}
\Level{A}
\SubArea{Infinite series}
\LANGen{
\Question{
We are  given an infinite geometric series:
\[ 4+\frac83+\frac{16}9+\frac{32}{27}+\dots\text{ .} \]
What is the value of its common ratio?}
{\( \frac23 \)}{\( \frac13 \)}{\( \frac43 \)}{\( \frac32 \)}{\( \frac34 \)}{}}
\LANGpl{
\Question{
Dany są ciągi geometryczne nieskończone 
\[ 4+\frac83+\frac{16}9+\frac{32}{27}+\dots\text{ .} \]
Jaka jest wartość ich stosunku?}
{\( \frac23 \)}{\( \frac13 \)}{\( \frac43 \)}{\( \frac32 \)}{\( \frac34 \)}{}}
\LANGcs{
\Question{
Je dána nekonečná geometrická řada:
\[ 4+\frac83+\frac{16}9+\frac{32}{27}+\dots\text{ .} \]
Její kvocient je roven:}
{\( \frac23 \)}{\( \frac13 \)}{\( \frac43 \)}{\( \frac32 \)}{\( \frac34 \)}{}}
\LANGsk{
\Question{
Je daný nekonečný geometrický rad:
\[ 4+\frac83+\frac{16}9+\frac{32}{27}+\dots\text{ .} \]
Jeho kvocient je rovný:}
{\( \frac23 \)}{\( \frac13 \)}{\( \frac43 \)}{\( \frac32 \)}{\( \frac34 \)}{}}
\LANGes{
\Question{
Dada la serie geométrica infinita:
\[ 4+\frac83+\frac{16}9+\frac{32}{27}+\dots\text{ .} \]
Su razón equivale a:}
{\( \frac23 \)}{\( \frac13 \)}{\( \frac43 \)}{\( \frac32 \)}{\( \frac34 \)}{}}
\End

\Data\ID{1003108706}
\Updated{2021-05-05 06:05:04}
\Level{A}
\SubArea{Infinite series}
\LANGen{
\Question{
We are given an infinite geometric series:
\[ \frac23-\frac49+\frac8{27}-\frac{16}{81}+\dots\text{ .} \]
What is the value of its common ratio?}
{\( -\frac23 \)}{\( \frac23 \)}{\( \frac29 \)}{\( -\frac29 \)}{\( \frac4{27} \)}{}}
\LANGpl{
\Question{
Dane są ciągi geometryczne nieskończone
\[ \frac23-\frac49+\frac8{27}-\frac{16}{81}+\dots\text{ .} \]
Jaki jest ich stosunek?}
{\( -\frac23 \)}{\( \frac23 \)}{\( \frac29 \)}{\( -\frac29 \)}{\( \frac4{27} \)}{}}
\LANGcs{
\Question{
Je dána nekonečná geometrická řada:
\[ \frac23-\frac49+\frac8{27}-\frac{16}{81}+\dots\text{ .} \]
Její kvocient je roven:}
{\( -\frac23 \)}{\( \frac23 \)}{\( \frac29 \)}{\( -\frac29 \)}{\( \frac4{27} \)}{}}
\LANGsk{
\Question{
Je daný nekonečný geometrický rad:
\[ \frac23-\frac49+\frac8{27}-\frac{16}{81}+\dots\text{ .} \]
Jeho kvocient je rovný:}
{\( -\frac23 \)}{\( \frac23 \)}{\( \frac29 \)}{\( -\frac29 \)}{\( \frac4{27} \)}{}}
\LANGes{
\Question{
Dada la serie geométrica infinita:
\[ \frac23-\frac49+\frac8{27}-\frac{16}{81}+\dots\text{ .} \]
Su razón equivale a:}
{\( -\frac23 \)}{\( \frac23 \)}{\( \frac29 \)}{\( -\frac29 \)}{\( \frac4{27} \)}{}}
\End

\Data\ID{1003108707}
\Updated{2021-05-05 06:05:52}
\Level{A}
\SubArea{Infinite series}
\LANGen{
\Question{
We are  given an infinite geometric series:
\[ \left(\sqrt5-\sqrt3\right)+\left(5-\sqrt{15}\right)+\left(5\sqrt5-5\sqrt{3}\right)+\dots\text{ .} \]
What is the value of its common ratio?}
{\( \sqrt5 \)}{\( \sqrt5-\sqrt3 \)}{\( \sqrt5-\sqrt3+5 \)}{\( \sqrt5+5 \)}{\( 5 \)}{}}
\LANGpl{
\Question{
Dane są szeregi geometryczne nieskończone:
\[ \left(\sqrt5-\sqrt3\right)+\left(5-\sqrt{15}\right)+\left(5\sqrt5-5\sqrt{3}\right)+\dots\text{ .} \]
Jaka jest wartość ich stosunku?}
{\( \sqrt5 \)}{\( \sqrt5-\sqrt3 \)}{\( \sqrt5-\sqrt3+5 \)}{\( \sqrt5+5 \)}{\( 5 \)}{}}
\LANGcs{
\Question{
Je dána nekonečná geometrická řada:
\[ \left(\sqrt5-\sqrt3\right)+\left(5-\sqrt{15}\right)+\left(5\sqrt5-5\sqrt{3}\right)+\dots\text{ .} \]
Její kvocient je roven:}
{\( \sqrt5 \)}{\( \sqrt5-\sqrt3 \)}{\( \sqrt5-\sqrt3+5 \)}{\( \sqrt5+5 \)}{\( 5 \)}{}}
\LANGsk{
\Question{
Je daný nekonečný geometrický rad:
\[ \left(\sqrt5-\sqrt3\right)+\left(5-\sqrt{15}\right)+\left(5\sqrt5-5\sqrt{3}\right)+\dots\text{ .} \]
Jeho kvocient je rovný:}
{\( \sqrt5 \)}{\( \sqrt5-\sqrt3 \)}{\( \sqrt5-\sqrt3+5 \)}{\( \sqrt5+5 \)}{\( 5 \)}{}}
\LANGes{
\Question{
Dada la serie geométrica infinita:
\[ \left(\sqrt5-\sqrt3\right)+\left(5-\sqrt{15}\right)+\left(5\sqrt5-5\sqrt{3}\right)+\dots\text{ .} \]
Su razón equivale a:}
{\( \sqrt5 \)}{\( \sqrt5-\sqrt3 \)}{\( \sqrt5-\sqrt3+5 \)}{\( \sqrt5+5 \)}{\( 5 \)}{}}
\End

\Data\ID{1003108708}
\Updated{2020-07-15 03:07:19}
\Level{A}
\SubArea{Infinite series}
\LANGen{
\Question{
We are  given an infinite geometric series:
\[ \sum\limits_{n=1}^{\infty} \left(-\frac12\right)^{n-1} \text{ .} \]
What is the value of its second term \( a_2 \)?}
{\( -\frac12 \)}{\( 1 \)}{\( \frac12 \)}{\( \frac14 \)}{\( -1 \)}{}}
\LANGpl{
\Question{
Dany jest szereg geometryczny nieskończony:
\[ \sum\limits_{n=1}^{\infty} \left(-\frac12\right)^{n-1} \text{ .} \]
Jaka jest wartość jego drugiego wyraz \( a_2 \)?}
{\( -\frac12 \)}{\( 1 \)}{\( \frac12 \)}{\( \frac14 \)}{\( -1 \)}{}}
\LANGcs{
\Question{
Je dána nekonečná geometrická řada:
\[ \sum\limits_{n=1}^{\infty} \left(-\frac12\right)^{n-1} \text{ .} \]
Její druhý člen  \( a_2 \) je roven:}
{\( -\frac12 \)}{\( 1 \)}{\( \frac12 \)}{\( \frac14 \)}{\( -1 \)}{}}
\LANGsk{
\Question{
Je daný nekonečný geometrický rad:
\[ \sum\limits_{n=1}^{\infty} \left(-\frac12\right)^{n-1} \text{ .} \]
Jeho druhý člen  \( a_2 \) je rovný:}
{\( -\frac12 \)}{\( 1 \)}{\( \frac12 \)}{\( \frac14 \)}{\( -1 \)}{}}
\LANGes{
\Question{
Dada la serie geométrica infinita:
\[ \sum\limits_{n=1}^{\infty} \left(-\frac12\right)^{n-1} \text{ .} \]
Su segundo término \( a_2 \) equivale a:}
{\( -\frac12 \)}{\( 1 \)}{\( \frac12 \)}{\( \frac14 \)}{\( -1 \)}{}}
\End

\Data\ID{1003108709}
\Updated{2020-07-15 03:07:19}
\Level{A}
\SubArea{Infinite series}
\LANGen{
\Question{
We are  given an infinite geometric series:
\[ \sum\limits_{n=1}^{\infty} \frac{\sqrt3}{2^{n-1}}\text{ .} \]
What is the value of its first term \( a_1 \)?}
{\( \sqrt3 \)}{\( \frac{\sqrt3}2 \)}{\( \frac{\sqrt3}4 \)}{\( \frac12 \)}{\( 3 \)}{}}
\LANGpl{
\Question{
Dany jest szereg geometryczny nieskończony:
\[ \sum\limits_{n=1}^{\infty} \frac{\sqrt3}{2^{n-1}}\text{ .} \]
Jaka jest wartość jego pierwszego wyrazu \( a_1 \)?}
{\( \sqrt3 \)}{\( \frac{\sqrt3}2 \)}{\( \frac{\sqrt3}4 \)}{\( \frac12 \)}{\( 3 \)}{}}
\LANGcs{
\Question{
Je dána nekonečná geometrická řada:
\[ \sum\limits_{n=1}^{\infty} \frac{\sqrt3}{2^{n-1}}\text{ .} \]
Její první člen  \( a_1 \) je roven:}
{\( \sqrt3 \)}{\( \frac{\sqrt3}2 \)}{\( \frac{\sqrt3}4 \)}{\( \frac12 \)}{\( 3 \)}{}}
\LANGsk{
\Question{
Je daný nekonečný geometrický rad:
\[ \sum\limits_{n=1}^{\infty} \frac{\sqrt3}{2^{n-1}}\text{ .} \]
Jej prvý člen  \( a_1 \) je rovný:}
{\( \sqrt3 \)}{\( \frac{\sqrt3}2 \)}{\( \frac{\sqrt3}4 \)}{\( \frac12 \)}{\( 3 \)}{}}
\LANGes{
\Question{
Dada la serie geométrica infinita:
\[ \sum\limits_{n=1}^{\infty} \frac{\sqrt3}{2^{n-1}}\text{ .} \]
Su primer término \( a_1 \) equivale a:}
{\( \sqrt3 \)}{\( \frac{\sqrt3}2 \)}{\( \frac{\sqrt3}4 \)}{\( \frac12 \)}{\( 3 \)}{}}
\End

\Data\ID{1003108710}
\Updated{2020-07-15 03:07:19}
\Level{A}
\SubArea{Infinite series}
\LANGen{
\Question{
The sum of the infinite geometric series
\[ \left(\sqrt5-2\right)+\left(\sqrt5-2\right)^2+\left(\sqrt5-2\right)^3+\dots \]
is equal to:}
{\( \frac{\sqrt5-1}4 \)}{\( \frac{\sqrt5}4 \)}{\( \frac{\sqrt5+1}4 \)}{\( \frac{\sqrt5+3}4 \)}{\( \frac{\sqrt5-1}2 \)}{}}
\LANGpl{
\Question{
Suma szeregów geometrycznych nieskończonych
\[ \left(\sqrt5-2\right)+\left(\sqrt5-2\right)^2+\left(\sqrt5-2\right)^3+\dots \]
jest równa:}
{\( \frac{\sqrt5-1}4 \)}{\( \frac{\sqrt5}4 \)}{\( \frac{\sqrt5+1}4 \)}{\( \frac{\sqrt5+3}4 \)}{\( \frac{\sqrt5-1}2 \)}{}}
\LANGcs{
\Question{
Součet nekonečné geometrické řady
\[ \left(\sqrt5-2\right)+\left(\sqrt5-2\right)^2+\left(\sqrt5-2\right)^3+\dots \]
je roven:}
{\( \frac{\sqrt5-1}4 \)}{\( \frac{\sqrt5}4 \)}{\( \frac{\sqrt5+1}4 \)}{\( \frac{\sqrt5+3}4 \)}{\( \frac{\sqrt5-1}2 \)}{}}
\LANGsk{
\Question{
Súčet nekonečného geometrického radu
\[ \left(\sqrt5-2\right)+\left(\sqrt5-2\right)^2+\left(\sqrt5-2\right)^3+\dots \]
je rovný:}
{\( \frac{\sqrt5-1}4 \)}{\( \frac{\sqrt5}4 \)}{\( \frac{\sqrt5+1}4 \)}{\( \frac{\sqrt5+3}4 \)}{\( \frac{\sqrt5-1}2 \)}{}}
\LANGes{
\Question{
La suma de la serie geométrica infinita
\[ \left(\sqrt5-2\right)+\left(\sqrt5-2\right)^2+\left(\sqrt5-2\right)^3+\dots \]
equivale a:}
{\( \frac{\sqrt5-1}4 \)}{\( \frac{\sqrt5}4 \)}{\( \frac{\sqrt5+1}4 \)}{\( \frac{\sqrt5+3}4 \)}{\( \frac{\sqrt5-1}2 \)}{}}
\End

\Data\ID{1003108711}
\Updated{2020-07-15 03:07:19}
\Level{A}
\SubArea{Infinite series}
\LANGen{
\Question{
The sum of the infinite geometric series
\[ 1-\frac34+\frac9{16}-\frac{27}{64}+\dots \]
is equal to:}
{\( \frac47 \)}{\( \frac14 \)}{\( 4 \)}{\( \frac74 \)}{\( \frac57 \)}{}}
\LANGpl{
\Question{
Suma szeregów geometrycznych nieskończonych 
\[ 1-\frac34+\frac9{16}-\frac{27}{64}+\dots \]
jest równa:}
{\( \frac47 \)}{\( \frac14 \)}{\( 4 \)}{\( \frac74 \)}{\( \frac57 \)}{}}
\LANGcs{
\Question{
Součet nekonečné geometrické řady
\[ 1-\frac34+\frac9{16}-\frac{27}{64}+\dots \]
je roven:}
{\( \frac47 \)}{\( \frac14 \)}{\( 4 \)}{\( \frac74 \)}{\( \frac57 \)}{}}
\LANGsk{
\Question{
Súčet nekonečného geometrického radu
\[ 1-\frac34+\frac9{16}-\frac{27}{64}+\dots \]
je rovný:}
{\( \frac47 \)}{\( \frac14 \)}{\( 4 \)}{\( \frac74 \)}{\( \frac57 \)}{}}
\LANGes{
\Question{
La suma de la serie geométrica infinita
\[ 1-\frac34+\frac9{16}-\frac{27}{64}+\dots \]
equivale a:}
{\( \frac47 \)}{\( \frac14 \)}{\( 4 \)}{\( \frac74 \)}{\( \frac57 \)}{}}
\End

\Data\ID{1003108712}
\Updated{2020-07-15 03:07:19}
\Level{A}
\SubArea{Infinite series}
\LANGen{
\Question{
The value of the expression
\[ \sum\limits_{n=1}^{\infty}(-1)^n\cdot\left(\frac23\right)^n \]
is equal to:}
{\( -\frac25 \)}{\( \frac25 \)}{\( -\frac23 \)}{\( \frac23 \)}{\( \frac52 \)}{}}
\LANGpl{
\Question{
Wartość wyrażenia
\[ \sum\limits_{n=1}^{\infty}(-1)^n\cdot\left(\frac23\right)^n \]
jest równa:}
{\( -\frac25 \)}{\( \frac25 \)}{\( -\frac23 \)}{\( \frac23 \)}{\( \frac52 \)}{}}
\LANGcs{
\Question{
Výraz
\[ \sum\limits_{n=1}^{\infty}(-1)^n\cdot\left(\frac23\right)^n \]
je roven:}
{\( -\frac25 \)}{\( \frac25 \)}{\( -\frac23 \)}{\( \frac23 \)}{\( \frac52 \)}{}}
\LANGsk{
\Question{
Výraz
\[ \sum\limits_{n=1}^{\infty}(-1)^n\cdot\left(\frac23\right)^n \]
je rovný:}
{\( -\frac25 \)}{\( \frac25 \)}{\( -\frac23 \)}{\( \frac23 \)}{\( \frac52 \)}{}}
\LANGes{
\Question{
El valor de la expresión
\[ \sum\limits_{n=1}^{\infty}(-1)^n\cdot\left(\frac23\right)^n \]
equivale a:}
{\( -\frac25 \)}{\( \frac25 \)}{\( -\frac23 \)}{\( \frac23 \)}{\( \frac52 \)}{}}
\End

\Data\ID{1003108713}
\Updated{2020-07-15 03:07:19}
\Level{A}
\SubArea{Infinite series}
\LANGen{
\Question{
The value of the expression
\[ \sum\limits_{n=1}^{\infty}\left(\frac13\right)^n \]
is equal to:}
{\( \frac12 \)}{\( 2 \)}{\( 1 \)}{\( \frac14 \)}{\( \frac23 \)}{}}
\LANGpl{
\Question{
Wartość wyrażenia \[ \sum\limits_{n=1}^{\infty}\left(\frac13\right)^n \]
jest równa:}
{\( \frac12 \)}{\( 2 \)}{\( 1 \)}{\( \frac14 \)}{\( \frac23 \)}{}}
\LANGcs{
\Question{
Výraz
\[ \sum\limits_{n=1}^{\infty}\left(\frac13\right)^n \]
je roven:}
{\( \frac12 \)}{\( 2 \)}{\( 1 \)}{\( \frac14 \)}{\( \frac23 \)}{}}
\LANGsk{
\Question{
Výraz
\[ \sum\limits_{n=1}^{\infty}\left(\frac13\right)^n \]
je rovný:}
{\( \frac12 \)}{\( 2 \)}{\( 1 \)}{\( \frac14 \)}{\( \frac23 \)}{}}
\LANGes{
\Question{
El valor de la expresión
\[ \sum\limits_{n=1}^{\infty}\left(\frac13\right)^n \]
equivale a:}
{\( \frac12 \)}{\( 2 \)}{\( 1 \)}{\( \frac14 \)}{\( \frac23 \)}{}}
\End

\Data\ID{9000062901}
\Updated{2020-07-15 03:07:37}
\Level{A}
\SubArea{Infinite series}
\LANGen{
\Question{
Find the sum of the following geometric series. 
\[ 
 -\frac{1}
{3} + \frac{1} 
{6} - \frac{1} 
{12} + \frac{1} 
{24}-\cdots 
\]}
{\(-\frac{2} 
{9}\)}{\(-\frac{2} 
{3}\)}{\(\frac{2}
{9}\)}{\(\infty \)}{}{}}
\LANGpl{
\Question{
Oblicz sumę następujących szeregów geometrycznych.
\[ 
 -\frac{1}
{3} + \frac{1} 
{6} - \frac{1} 
{12} + \frac{1} 
{24}-\cdots 
\]}
{\(-\frac{2} 
{9}\)}{\(-\frac{2} 
{3}\)}{\(\frac{2}
{9}\)}{\(\infty \)}{}{}}
\LANGcs{
\Question{
Určete součet geometrické řady.
\[-\frac{1} 
{3} + \frac{1} 
{6} - \frac{1} 
{12} + \frac{1} 
{24}-\cdots \]}
{\(-\frac{2} 
{9}\)}{\(-\frac{2} 
{3}\)}{\(\frac{2}
{9}\)}{\(\infty\)}{}{}}
\LANGsk{
\Question{
Určte súčet geometrického radu.
\[-\frac{1} 
{3} + \frac{1} 
{6} - \frac{1} 
{12} + \frac{1} 
{24}-\cdots \]}
{\(-\frac{2} 
{9}\)}{\(-\frac{2} 
{3}\)}{\(\frac{2}
{9}\)}{\(\infty\)}{}{}}
\LANGes{
\Question{
Halla la suma de la siguiente serie geométrica:
\[ 
 -\frac{1}
{3} + \frac{1} 
{6} - \frac{1} 
{12} + \frac{1} 
{24}-\cdots 
\]}
{\(-\frac{2} 
{9}\)}{\(-\frac{2} 
{3}\)}{\(\frac{2}
{9}\)}{\(\infty \)}{}{}}
\End

\Data\ID{9000062902}
\Updated{2020-07-15 03:07:37}
\Level{A}
\SubArea{Infinite series}
\LANGen{
\Question{
Find the sum of the following geometric series. 
\[ 
 1 + \frac{3} 
{2} + \frac{9} 
{4} + \frac{27} 
 {8} + \frac{81} 
{16}+\cdots 
\]}
{\(\infty \)}{\(- 2\)}{\(2\)}{\(\frac{2}
{5}\)}{}{}}
\LANGpl{
\Question{
Oblicz sumę następujących szeregów geometrycznych.
\[ 
 1 + \frac{3} 
{2} + \frac{9} 
{4} + \frac{27} 
 {8} + \frac{81} 
{16}+\cdots 
\]}
{\(\infty \)}{\(- 2\)}{\(2\)}{\(\frac{2}
{5}\)}{}{}}
\LANGcs{
\Question{
Určete součet geometrické řady.
\[1 + \frac{3} 
{2} + \frac{9} 
{4} + \frac{27} 
 {8} + \frac{81} 
{16}+\cdots \]}
{\(\infty\)}{\(- 2\)}{\(2\)}{\(\frac{2}{5}\)}{}{}}
\LANGsk{
\Question{
Určte súčet geometrického radu.
\[1 + \frac{3} 
{2} + \frac{9} 
{4} + \frac{27} 
 {8} + \frac{81} 
{16}+\cdots \]}
{\(\infty\)}{\(- 2\)}{\(2\)}{\(\frac{2}{5}\)}{}{}}
\LANGes{
\Question{
Halla la suma de la siguiente serie geométrica:
\[ 
 1 + \frac{3} 
{2} + \frac{9} 
{4} + \frac{27} 
 {8} + \frac{81} 
{16}+\cdots 
\]}
{\(\infty \)}{\(- 2\)}{\(2\)}{\(\frac{2}
{5}\)}{}{}}
\End

\Data\ID{9000063401}
\Updated{2021-05-05 06:05:02}
\Level{A}
\SubArea{Infinite series}
\LANGen{
\Question{
Identify the quotient of the geometric series
\(\sum _{n=1}^{\infty } \frac{1}
{2^{n-3}} \).}
{\(\frac{1}
{2}\)}{\(2\)}{\(1\)}{\(\frac{1}
{8}\)}{}{}}
\LANGpl{
\Question{
Oblicz iloraz szeregów geometrycznych
\(\sum _{n=1}^{\infty } \frac{1}
{2^{n-3}} \).}
{\(\frac{1}
{2}\)}{\(2\)}{\(1\)}{\(\frac{1}
{8}\)}{}{}}
\LANGcs{
\Question{
Je dána nekonečná geometrická řada
\(\sum _{n=1}^{\infty } \frac{1}
{2^{n-3}} \). Její
kvocient \(q\) 
je roven:}
{\(\frac{1}
{2}\)}{\(2\)}{\(1\)}{\(\frac{1}
{8}\)}{}{}}
\LANGsk{
\Question{
Je daný nekonečný geometrický rad
\(\sum _{n=1}^{\infty } \frac{1}
{2^{n-3}} \). Jeho
kvocient \(q\) 
je rovný:}
{\(\frac{1}
{2}\)}{\(2\)}{\(1\)}{\(\frac{1}
{8}\)}{}{}}
\LANGes{
\Question{
Dada la serie geométrica infinita:
\(\sum _{n=1}^{\infty } \frac{1}
{2^{n-3}} \). Su razón \(q\) 
equivale a:}
{\(\frac{1}
{2}\)}{\(2\)}{\(1\)}{\(\frac{1}
{8}\)}{}{}}
\End

\Data\ID{9000063402}
\Updated{2021-05-05 06:05:35}
\Level{A}
\SubArea{Infinite series}
\LANGen{
\Question{
Identify the quotient of the geometric series
\(\sum _{n=1}^{\infty }3^{2-n}\).}
{\(\frac{1}
{3}\)}{\(1\)}{\(\frac{1}
{9}\)}{\(-\frac{1} 
{9}\)}{}{}}
\LANGpl{
\Question{
Oblicz iloraz szeregu geometrycznego
\(\sum _{n=1}^{\infty }3^{2-n}\).}
{\(\frac{1}
{3}\)}{\(1\)}{\(\frac{1}
{9}\)}{\(-\frac{1} 
{9}\)}{}{}}
\LANGcs{
\Question{
Je dána nekonečná geometrická řada
\(\sum _{n=1}^{\infty }3^{2-n}\). Její
kvocient \(q\) 
je roven:}
{\(\frac{1}
{3}\)}{\(1\)}{\(\frac{1}
{9}\)}{\(-\frac{1} 
{9}\)}{}{}}
\LANGsk{
\Question{
Je daný nekonečný geometrický rad
\(\sum _{n=1}^{\infty }3^{2-n}\). Jeho
kvocient \(q\) 
je rovný:}
{\(\frac{1}
{3}\)}{\(1\)}{\(\frac{1}
{9}\)}{\(-\frac{1} 
{9}\)}{}{}}
\LANGes{
\Question{
Dada la serie geométrica infinita:
\(\sum _{n=1}^{\infty }3^{2-n}\). Su razón \(q\) equivale a:}
{\(\frac{1}
{3}\)}{\(1\)}{\(\frac{1}
{9}\)}{\(-\frac{1} 
{9}\)}{}{}}
\End

\Data\ID{9000063403}
\Updated{2020-07-15 03:07:37}
\Level{A}
\SubArea{Infinite series}
\LANGen{
\Question{
Evaluate the following infinite product. 
\[ 
 2\cdot \sqrt{2}\cdot \root{4}\of{2}\cdot \root{8}\of{2}\cdot \cdots 
\]}
{\(4\)}{\(1\)}{\(2\)}{\(8\)}{}{}}
\LANGpl{
\Question{
Oblicz następujący iloczyn nieskończony.
\[ 
 2\cdot \sqrt{2}\cdot \root{4}\of{2}\cdot \root{8}\of{2}\cdot \cdots 
\]}
{\(4\)}{\(1\)}{\(2\)}{\(8\)}{}{}}
\LANGcs{
\Question{
Výraz \(2\cdot \sqrt{2}\cdot \root{4}\of{2}\cdot \root{8}\of{2}\cdot \cdots \) 
je roven:}
{\(4\)}{\(1\)}{\(2\)}{\(8\)}{}{}}
\LANGsk{
\Question{
Výraz \(2\cdot \sqrt{2}\cdot \root{4}\of{2}\cdot \root{8}\of{2}\cdot \cdots \) 
je rovný:}
{\(4\)}{\(1\)}{\(2\)}{\(8\)}{}{}}
\LANGes{
\Question{
La expresión
\[ 
 2\cdot \sqrt{2}\cdot \root{4}\of{2}\cdot \root{8}\of{2}\cdot \cdots 
\] equivale a:}
{\(4\)}{\(1\)}{\(2\)}{\(8\)}{}{}}
\End

\Data\ID{9000063404}
\Updated{2020-07-15 03:07:37}
\Level{A}
\SubArea{Infinite series}
\LANGen{
\Question{
Evaluate the following infinite sum. 
\[ 
 \frac{5} 
{2} + \frac{5} 
{8} + \frac{5} 
{32} + \frac{5} 
{128}+\cdots 
\]}
{\(\frac{10}
 {3} \)}{\(5\)}{\(4\)}{\(\frac{5}
{2}\)}{}{}}
\LANGpl{
\Question{
Oblicz następującą sumę nieskończoną.
\[ 
 \frac{5} 
{2} + \frac{5} 
{8} + \frac{5} 
{32} + \frac{5} 
{128}+\cdots 
\]}
{\(\frac{10}
 {3} \)}{\(5\)}{\(4\)}{\(\frac{5}
{2}\)}{}{}}
\LANGcs{
\Question{
Výraz \(\frac{5}
{2} + \frac{5} 
{8} + \frac{5} 
{32} + \frac{5} 
{128}+\cdots \) 
je roven:}
{\(\frac{10}
 {3} \)}{\(5\)}{\(4\)}{\(\frac{5}
{2}\)}{}{}}
\LANGsk{
\Question{
Výraz \(\frac{5}
{2} + \frac{5} 
{8} + \frac{5} 
{32} + \frac{5} 
{128}+\cdots \) 
je rovný:}
{\(\frac{10}
 {3} \)}{\(5\)}{\(4\)}{\(\frac{5}
{2}\)}{}{}}
\LANGes{
\Question{
La expresión
\[ 
 \frac{5} 
{2} + \frac{5} 
{8} + \frac{5} 
{32} + \frac{5} 
{128}+\cdots 
\] equivale a:}
{\(\frac{10}
 {3} \)}{\(5\)}{\(4\)}{\(\frac{5}
{2}\)}{}{}}
\End

\Data\ID{9000063405}
\Updated{2020-07-15 03:07:37}
\Level{A}
\SubArea{Infinite series}
\LANGen{
\Question{
Evaluate the following infinite sum. 
\[ 
 -\frac{2}
{3} + \frac{1} 
{6} -\frac{2} 
{6} + \frac{1} 
{12} - \frac{2} 
{12} + \frac{1} 
{24}+\cdots 
\]}
{\(- 1\)}{\(-\frac{4} 
{3}\)}{\(\frac{1}
{3}\)}{\(\frac{3}
{2}\)}{}{}}
\LANGpl{
\Question{
Oblicz następująca sumę nieskończoną.
\[ 
 -\frac{2}
{3} + \frac{1} 
{6} -\frac{2} 
{6} + \frac{1} 
{12} - \frac{2} 
{12} + \frac{1} 
{24}+\cdots 
\]}
{\(- 1\)}{\(-\frac{4} 
{3}\)}{\(\frac{1}
{3}\)}{\(\frac{3}
{2}\)}{}{}}
\LANGcs{
\Question{
Výraz \(-\frac{2} 
{3} + \frac{1} 
{6} -\frac{2} 
{6} + \frac{1} 
{12} - \frac{2} 
{12} + \frac{1} 
{24}+\cdots \) 
je roven:}
{\(- 1\)}{\(-\frac{4} 
{3}\)}{\(\frac{1}
{3}\)}{\(\frac{3}
{2}\)}{}{}}
\LANGsk{
\Question{
Výraz \(-\frac{2} 
{3} + \frac{1} 
{6} -\frac{2} 
{6} + \frac{1} 
{12} - \frac{2} 
{12} + \frac{1} 
{24}+\cdots \) 
je rovný:}
{\(- 1\)}{\(-\frac{4} 
{3}\)}{\(\frac{1}
{3}\)}{\(\frac{3}
{2}\)}{}{}}
\LANGes{
\Question{
La expresión
\[ 
 -\frac{2}
{3} + \frac{1} 
{6} -\frac{2} 
{6} + \frac{1} 
{12} - \frac{2} 
{12} + \frac{1} 
{24}+\cdots 
\] equivale a:}
{\(- 1\)}{\(-\frac{4} 
{3}\)}{\(\frac{1}
{3}\)}{\(\frac{3}
{2}\)}{}{}}
\End

\Data\ID{9000063406}
\Updated{2020-07-15 03:07:37}
\Level{A}
\SubArea{Infinite series}
\LANGen{
\Question{
Evaluate the following infinite sum. 
\[ 
 \sum _{n=1}^{\infty }\left (-\frac{1}
{2}\right )^{n+2}
\]}
{\(- \frac{1} 
{12}\)}{\(-\frac{1} 
{8}\)}{\(\frac{1}
{2}\)}{\(1\)}{}{}}
\LANGpl{
\Question{
Oblicz następującą sumę nieskończoną.
\[ 
 \sum _{n=1}^{\infty }\left (-\frac{1}
{2}\right )^{n+2}
\]}
{\(- \frac{1} 
{12}\)}{\(-\frac{1} 
{8}\)}{\(\frac{1}
{2}\)}{\(1\)}{}{}}
\LANGcs{
\Question{
Výraz \(\sum _{n=1}^{\infty }\left (-\frac{1}
{2}\right )^{n+2}\) 
je roven:}
{\(- \frac{1} 
{12}\)}{\(-\frac{1} 
{8}\)}{\(\frac{1}
{2}\)}{\(1\)}{}{}}
\LANGsk{
\Question{
Výraz \(\sum _{n=1}^{\infty }\left (-\frac{1}
{2}\right )^{n+2}\) 
je rovný:}
{\(- \frac{1} 
{12}\)}{\(-\frac{1} 
{8}\)}{\(\frac{1}
{2}\)}{\(1\)}{}{}}
\LANGes{
\Question{
La expresión
\[ 
 \sum _{n=1}^{\infty }\left (-\frac{1}
{2}\right )^{n+2}
\] equivale a:}
{\(- \frac{1} 
{12}\)}{\(-\frac{1} 
{8}\)}{\(\frac{1}
{2}\)}{\(1\)}{}{}}
\End

\Data\ID{9000073404}
\Updated{2020-07-15 03:07:37}
\Level{A}
\SubArea{Infinite series}
\LANGen{
\Question{
Find the sum of the following infinite series. 
\[ 
 \sqrt{2} - 2 + \sqrt{8} - 4 + \sqrt{32} - 8+\cdots 
\]}
{The sum does not exist.}{\(\frac{\sqrt{2}}
{1+\sqrt{2}}\)}{\(\frac{\sqrt{2}}
{1-\sqrt{2}}\)}{\(\sqrt{2} - 2\)}{}{}}
\LANGpl{
\Question{
Oblicz sumę następujących szeregów nieskończonych.
\[ 
 \sqrt{2} - 2 + \sqrt{8} - 4 + \sqrt{32} - 8+\cdots 
\]}
{Szereg jest rozbieżny.}{\(\frac{\sqrt{2}}
{1+\sqrt{2}}\)}{\(\frac{\sqrt{2}}
{1-\sqrt{2}}\)}{\(\sqrt{2} - 2\)}{}{}}
\LANGcs{
\Question{
Určete, zda nekonečná řada \(\sqrt{2} - 2 + \sqrt{8} - 4 + \sqrt{32} - 8+\cdots \) 
konverguje nebo diverguje. V případě, že konverguje, určete její součet.}
{Řada je divergentní.}{\(\frac{\sqrt{2}}
{1+\sqrt{2}}\)}{\(\frac{\sqrt{2}}
{1-\sqrt{2}}\)}{\(\sqrt{2} - 2\)}{}{}}
\LANGsk{
\Question{
Určte, či nekonečný rad \(\sqrt{2} - 2 + \sqrt{8} - 4 + \sqrt{32} - 8+\cdots \) 
konverguje alebo diverguje. V prípade, že konverguje, určte jeho súčet.}
{Rad je divergentný.}{\(\frac{\sqrt{2}}
{1+\sqrt{2}}\)}{\(\frac{\sqrt{2}}
{1-\sqrt{2}}\)}{\(\sqrt{2} - 2\)}{}{}}
\LANGes{
\Question{
Halla la suma de la siguiente serie infinita:
\[ 
 \sqrt{2} - 2 + \sqrt{8} - 4 + \sqrt{32} - 8+\cdots 
\]}
{La suma no existe.}{\(\frac{\sqrt{2}}
{1+\sqrt{2}}\)}{\(\frac{\sqrt{2}}
{1-\sqrt{2}}\)}{\(\sqrt{2} - 2\)}{}{}}
\End

\Data\ID{9000073405}
\Updated{2020-07-15 03:07:37}
\Level{A}
\SubArea{Infinite series}
\LANGen{
\Question{
Find the sum of the following infinite series. 
\[ 
 \sqrt{2} - 1 + \frac{\sqrt{2}} 
 {2} -\frac{1} 
{2} + \frac{\sqrt{2}} 
 {4} -\frac{1} 
{4}+\cdots 
\]}
{\(2\sqrt{2} - 2\)}{\(\sqrt{2} - 1\)}{\(2\sqrt{2} + 2\)}{\(\infty \)}{}{}}
\LANGpl{
\Question{
Oblicz sumę następujących szeregów nieskończonych.
\[ 
 \sqrt{2} - 1 + \frac{\sqrt{2}} 
 {2} -\frac{1} 
{2} + \frac{\sqrt{2}} 
 {4} -\frac{1} 
{4}+\cdots 
\]}
{\(2\sqrt{2} - 2\)}{\(\sqrt{2} - 1\)}{\(2\sqrt{2} + 2\)}{\(\infty \)}{}{}}
\LANGcs{
\Question{
Určete, zda nekonečná řada \(\sqrt{2} - 1 + \frac{\sqrt{2}} 
 {2} -\frac{1} 
{2} + \frac{\sqrt{2}} 
 {4} -\frac{1} 
{4}+\cdots \) 
konverguje nebo diverguje. V případě, že konverguje, určete její součet.}
{\(2\sqrt{2} - 2\)}{\(\sqrt{2} - 1\)}{\(2\sqrt{2} + 2\)}{$\infty$}{}{}}
\LANGsk{
\Question{
Určte, či nekonečný rad \(\sqrt{2} - 1 + \frac{\sqrt{2}} 
 {2} -\frac{1} 
{2} + \frac{\sqrt{2}} 
 {4} -\frac{1} 
{4}+\cdots \) 
konverguje alebo diverguje. V prípade, že konverguje, určte jeho súčet.}
{\(2\sqrt{2} - 2\)}{\(\sqrt{2} - 1\)}{\(2\sqrt{2} + 2\)}{$\infty$}{}{}}
\LANGes{
\Question{
Halla la suma de la siguiente serie infinita:
\[ 
 \sqrt{2} - 1 + \frac{\sqrt{2}} 
 {2} -\frac{1} 
{2} + \frac{\sqrt{2}} 
 {4} -\frac{1} 
{4}+\cdots 
\]}
{\(2\sqrt{2} - 2\)}{\(\sqrt{2} - 1\)}{\(2\sqrt{2} + 2\)}{\(\infty \)}{}{}}
\End

\Data\ID{9000073406}
\Updated{2021-01-08 03:01:50}
\Level{A}
\SubArea{Infinite series}
\LANGen{
\Question{
Find the sum of the following infinite series. 
\[ 
 \sum _{n=1}^{\infty }\left (\frac{\sqrt{2} - 1}
 {\sqrt{2}} \right )^{n-1}
\]}
{\(\sqrt{2}\)}{\(\frac{\sqrt{2}+1}
 {\sqrt{2}} \)}{\(\frac{\sqrt{2}}
 {2} \)}{Series diverges.}{}{}}
\LANGpl{
\Question{
Oblicz sumę następujących szeregów nieskończonych.
\[ 
 \sum _{n=1}^{\infty }\left (\frac{\sqrt{2} - 1}
 {\sqrt{2}} \right )^{n-1}
\]}
{\(\sqrt{2}\)}{\(\frac{\sqrt{2}+1}
 {\sqrt{2}} \)}{\(\frac{\sqrt{2}}
 {2} \)}{Szereg jest rozbieżny.}{}{}}
\LANGcs{
\Question{
Určete, zda nekonečná řada \(\sum _{n=1}^{\infty }\left (\frac{\sqrt{2}-1}
 {\sqrt{2}} \right )^{n-1}\) 
konverguje nebo diverguje. V případě, že konverguje, určete její součet.}
{\(\sqrt{2}\)}{\(\frac{\sqrt{2}+1}
 {\sqrt{2}} \)}{\(\frac{\sqrt{2}}
 {2} \)}{Řada je divergentní.}{}{}}
\LANGsk{
\Question{
Určte, či nekonečný rad \(\sum _{n=1}^{\infty }\left (\frac{\sqrt{2}-1}
 {\sqrt{2}} \right )^{n-1}\) 
konverguje alebo diverguje. V prípade, že konverguje, určte jeho súčet.}
{\(\sqrt{2}\)}{\(\frac{\sqrt{2}+1}
 {\sqrt{2}} \)}{\(\frac{\sqrt{2}}
 {2} \)}{Rad je divergentný.}{}{}}
\LANGes{
\Question{
Halla la suma de la siguiente serie infinita:
\[ 
 \sum _{n=1}^{\infty }\left (\frac{\sqrt{2} - 1}
 {\sqrt{2}} \right )^{n-1}
\]}
{\(\sqrt{2}\)}{\(\frac{\sqrt{2}+1}
 {\sqrt{2}} \)}{\(\frac{\sqrt{2}}
 {2} \)}{Es divergente.}{}{}}
\End

\Data\ID{1003108801}
\Updated{2021-05-05 06:05:38}
\Level{B}
\SubArea{Infinite series}
\LANGen{
\Question{
Express the repeating decimal \( 2.\overline{4} \) as a fraction in lowest terms.}
{\( \frac{22}9 \)}{\( \frac49 \)}{\( \frac{22}7 \)}{\( \frac{12}5 \)}{\( \frac9{22} \)}{}}
\LANGpl{
\Question{
Wyraź powtarzający się ułamek dziesiętny \( 2{,}\overline{4} \) jako ułamek najmniejszych wyrazów.}
{\( \frac{22}9 \)}{\( \frac49 \)}{\( \frac{22}7 \)}{\( \frac{12}5 \)}{\( \frac9{22} \)}{}}
\LANGcs{
\Question{
Periodické číslo  \( 2{,}\overline{4} \) zapište zlomkem v základním tvaru.}
{\( \frac{22}9 \)}{\( \frac49 \)}{\( \frac{22}7 \)}{\( \frac{12}5 \)}{\( \frac9{22} \)}{}}
\LANGsk{
\Question{
Periodické číslo  \( 2{,}\overline{4} \) zapíšte v tvare zlomku v základnom tvare.}
{\( \frac{22}9 \)}{\( \frac49 \)}{\( \frac{22}7 \)}{\( \frac{12}5 \)}{\( \frac9{22} \)}{}}
\LANGes{
\Question{
Expresa el decimal repetido \( 2.\overline{4} \) como fracción irreducible.}
{\( \frac{22}9 \)}{\( \frac49 \)}{\( \frac{22}7 \)}{\( \frac{12}5 \)}{\( \frac9{22} \)}{}}
\End

\Data\ID{1003108802}
\Updated{2021-05-05 06:05:12}
\Level{B}
\SubArea{Infinite series}
\LANGen{
\Question{
Express the repeating decimal \( 0.5\overline{83} \) as a fraction in lowest terms.}
{\( \frac{289}{495} \)}{\( \frac{83}{990} \)}{\( \frac{495}{289} \)}{\( \frac{298}{495} \)}{\( \frac{583}{1\,000} \)}{}}
\LANGpl{
\Question{
Wyraź powtarzający się ułamek dziesiętny \( 0{,}5\overline{83} \) jako ułamek najmniejszych wyrazów.}
{\( \frac{289}{495} \)}{\( \frac{83}{990} \)}{\( \frac{495}{289} \)}{\( \frac{298}{495} \)}{\( \frac{583}{1\,000} \)}{}}
\LANGcs{
\Question{
Periodické číslo  \( 0{,}5\overline{83} \) zapište zlomkem v základním tvaru.}
{\( \frac{289}{495} \)}{\( \frac{83}{990} \)}{\( \frac{495}{289} \)}{\( \frac{298}{495} \)}{\( \frac{583}{1\,000} \)}{}}
\LANGsk{
\Question{
Periodické číslo  \( 0{,}5\overline{83} \) zapíšte v tvare zlomku v základnom tvare.}
{\( \frac{289}{495} \)}{\( \frac{83}{990} \)}{\( \frac{495}{289} \)}{\( \frac{298}{495} \)}{\( \frac{583}{1\,000} \)}{}}
\LANGes{
\Question{
Expresa el decimal repetido \( 0.5\overline{83} \) como fracción irreducible.}
{\( \frac{289}{495} \)}{\( \frac{83}{990} \)}{\( \frac{495}{289} \)}{\( \frac{298}{495} \)}{\( \frac{583}{1\,000} \)}{}}
\End

\Data\ID{1003108803}
\Updated{2021-05-05 06:05:48}
\Level{B}
\SubArea{Infinite series}
\LANGen{
\Question{
What is the sum of the repeating numbers \( 0.\overline{27} \) and \( 0.4\overline{63} \)?}
{\( \frac{81}{110} \)}{\( \frac{110}{81} \)}{\( \frac{80}{111} \)}{\( \frac{79}{110} \)}{\( \frac{81}{100} \)}{}}
\LANGpl{
\Question{
Jaka jest suma powtarzających się liczb \( 0{,}\overline{27} \) i \( 0{,}4\overline{63} \)?}
{\( \frac{81}{110} \)}{\( \frac{110}{81} \)}{\( \frac{80}{111} \)}{\( \frac{79}{110} \)}{\( \frac{81}{100} \)}{}}
\LANGcs{
\Question{
Součet periodických čísel  \( 0{,}\overline{27} \)  a  \( 0{,}4\overline{63} \) je roven:}
{\( \frac{81}{110} \)}{\( \frac{110}{81} \)}{\( \frac{80}{111} \)}{\( \frac{79}{110} \)}{\( \frac{81}{100} \)}{}}
\LANGsk{
\Question{
Súčet periodických čísel  \( 0{,}\overline{27} \)  a  \( 0{,}4\overline{63} \) je rovný:}
{\( \frac{81}{110} \)}{\( \frac{110}{81} \)}{\( \frac{80}{111} \)}{\( \frac{79}{110} \)}{\( \frac{81}{100} \)}{}}
\LANGes{
\Question{
Halla la suma de los números periódicos \( 0.\overline{27} \) y \( 0.4\overline{63} \).}
{\( \frac{81}{110} \)}{\( \frac{110}{81} \)}{\( \frac{80}{111} \)}{\( \frac{79}{110} \)}{\( \frac{81}{100} \)}{}}
\End

\Data\ID{1003108804}
\Updated{2020-07-15 03:07:19}
\Level{B}
\SubArea{Infinite series}
\LANGen{
\Question{
What is the expression \( \frac{0.4\overline6}{0.\overline{63}} \) equal to?}
{\( \frac{11}{15} \)}{\( \frac{15}{11} \)}{\( \frac{11}{14} \)}{\( \frac{14}{11} \)}{\( \frac{13}{15} \)}{}}
\LANGpl{
\Question{
Wyrażenie \( \frac{0{,}4\overline6}{0{,}\overline{63}} \) równe jest:}
{\( \frac{11}{15} \)}{\( \frac{15}{11} \)}{\( \frac{11}{14} \)}{\( \frac{14}{11} \)}{\( \frac{13}{15} \)}{}}
\LANGcs{
\Question{
Výraz  \( \frac{0{,}4\overline6}{0{,}\overline{63}} \) je roven:}
{\( \frac{11}{15} \)}{\( \frac{15}{11} \)}{\( \frac{11}{14} \)}{\( \frac{14}{11} \)}{\( \frac{13}{15} \)}{}}
\LANGsk{
\Question{
Výraz  \( \frac{0{,}4\overline6}{0{,}\overline{63}} \) je rovný:}
{\( \frac{11}{15} \)}{\( \frac{15}{11} \)}{\( \frac{11}{14} \)}{\( \frac{14}{11} \)}{\( \frac{13}{15} \)}{}}
\LANGes{
\Question{
La expresión \( \frac{0.4\overline6}{0.\overline{63}} \) equivale a:}
{\( \frac{11}{15} \)}{\( \frac{15}{11} \)}{\( \frac{11}{14} \)}{\( \frac{14}{11} \)}{\( \frac{13}{15} \)}{}}
\End

\Data\ID{1003108805}
\Updated{2020-07-15 03:07:19}
\Level{B}
\SubArea{Infinite series}
\LANGen{
\Question{
We are given the equation 
\[ \sum\limits_{n=0}^{\infty}\left(\frac2x\right)^n=\frac{4x-3}{3x-4}  \]
with the unknown \( x \) being a real number. What is the set of all its solutions?}
{\( \{ 6 \} \)}{\( \{ 1 \} \)}{\( \{ 6;1 \} \)}{\( \{\ \} \)}{\( \{-6;-1\} \)}{}}
\LANGpl{
\Question{
Dane jest równanie
\[ \sum\limits_{n=0}^{\infty}\left(\frac2x\right)^n=\frac{4x-3}{3x-4}  \]
z niewiadomą  \( x \), która jest liczbą rzeczywistą. Jaki jest zbiór jego wszystkich rozwiązań?}
{\( \{ 6 \} \)}{\( \{ 1 \} \)}{\( \{ 6;1 \} \)}{\( \{\ \} \)}{\( \{-6;-1\} \)}{}}
\LANGcs{
\Question{
Je dána rovnice 
\[ \sum\limits_{n=0}^{\infty}\left(\frac2x\right)^n=\frac{4x-3}{3x-4}  \]
s reálnou neznámou  \( x \). Jaká je množina všech řešení této rovnice?}
{\( \{ 6 \} \)}{\( \{ 1 \} \)}{\( \{ 6;1 \} \)}{\( \{\ \} \)}{\( \{-6;-1\} \)}{}}
\LANGsk{
\Question{
Je daná rovnica 
\[ \sum\limits_{n=0}^{\infty}\left(\frac2x\right)^n=\frac{4x-3}{3x-4}  \]
s reálnou neznámou  \( x \). Aká je množina všetkých riešení tejto rovnice?}
{\( \{ 6 \} \)}{\( \{ 1 \} \)}{\( \{ 6;1 \} \)}{\( \{\ \} \)}{\( \{-6;-1\} \)}{}}
\LANGes{
\Question{
Dada la ecuación
\[ \sum\limits_{n=0}^{\infty}\left(\frac2x\right)^n=\frac{4x-3}{3x-4}  \]
con una incógnita  real \( x \). Calcula el conjunto de todas sus soluciones.}
{\( \{ 6 \} \)}{\( \{ 1 \} \)}{\( \{ 6;1 \} \)}{\( \{\ \} \)}{\( \{-6;-1\} \)}{}}
\End

\Data\ID{1003108806}
\Updated{2020-07-15 03:07:19}
\Level{B}
\SubArea{Infinite series}
\LANGen{
\Question{
We are given the equation
\[ \sum\limits_{n=1}^{\infty}(x+2)^{2n} =\frac13 \]
with the unknown \( x \) being a real number. What is the set of all its solutions?}
{\( \left\{ -\frac32;-\frac52 \right\} \)}{\( \left\{ -\frac32\right\} \)}{\( \left\{ -\frac52 \right\} \)}{\( \left\{\ \right\} \)}{\( \left\{ \frac32;\frac52 \right\} \)}{}}
\LANGpl{
\Question{
Dane jest równanie
\[ \sum\limits_{n=1}^{\infty}(x+2)^{2n} =\frac13 \]
z niewiadomą \( x \), która jest liczbą rzeczywistą. Jaki jest zbiór jego wszystkich rozwiązań?}
{\( \left\{ -\frac32;-\frac52 \right\} \)}{\( \left\{ -\frac32\right\} \)}{\( \left\{ -\frac52 \right\} \)}{\( \left\{\ \right\} \)}{\( \left\{ \frac32;\frac52 \right\} \)}{}}
\LANGcs{
\Question{
Je dána rovnice
\[ \sum\limits_{n=1}^{\infty}(x+2)^{2n} =\frac13 \]
s reálnou neznámou  \( x \). Jaká je množina všech řešení této rovnice?}
{\( \left\{ -\frac32;-\frac52 \right\} \)}{\( \left\{ -\frac32\right\} \)}{\( \left\{ -\frac52 \right\} \)}{\( \left\{\ \right\} \)}{\( \left\{ \frac32;\frac52 \right\} \)}{}}
\LANGsk{
\Question{
Je daná rovnica
\[ \sum\limits_{n=1}^{\infty}(x+2)^{2n} =\frac13 \]
s reálnou neznámou  \( x \). Aká je množina všetkých riešení tejto rovnice?}
{\( \left\{ -\frac32;-\frac52 \right\} \)}{\( \left\{ -\frac32\right\} \)}{\( \left\{ -\frac52 \right\} \)}{\( \left\{\ \right\} \)}{\( \left\{ \frac32;\frac52 \right\} \)}{}}
\LANGes{
\Question{
Dada la ecuación
\[ \sum\limits_{n=1}^{\infty}(x+2)^{2n} =\frac13 \]
con una incógnita real \( x \). Calcula el conjunto de todas sus soluciones.}
{\( \left\{ -\frac32;-\frac52 \right\} \)}{\( \left\{ -\frac32\right\} \)}{\( \left\{ -\frac52 \right\} \)}{\( \left\{\ \right\} \)}{\( \left\{ \frac32;\frac52 \right\} \)}{}}
\End

\Data\ID{1003108807}
\Updated{2020-07-15 03:07:19}
\Level{B}
\SubArea{Infinite series}
\LANGen{
\Question{
We are given the equation 
\[\log x+\log\sqrt x+\log\sqrt[4]x+\log\sqrt[8]x+\dots=2\]
with the unknown \( x \) being a real number. What is the set of all its solutions?}
{\( \{ 10 \} \)}{\( \{ 1 \} \)}{\( \{ 100 \} \)}{\( \{ \ \} \)}{\( \left\{ \frac1{10} \right\} \)}{}}
\LANGpl{
\Question{
Dane jest równanie
\[\log x+\log\sqrt x+\log\sqrt[4]x+\log\sqrt[8]x+\dots=2\]
z niewiadomą \( x \), która jest liczbą rzeczywistą.  Wskaż zbiór jego wszystkich rozwiązań.}
{\( \{ 10 \} \)}{\( \{ 1 \} \)}{\( \{ 100 \} \)}{\( \{ \ \} \)}{\( \left\{ \frac1{10} \right\} \)}{}}
\LANGcs{
\Question{
Je dána rovnice 
\[\log x+\log\sqrt x+\log\sqrt[4]x+\log\sqrt[8]x+\dots=2\]
s reálnou neznámou  \( x \). Jaká je množina všech řešení této rovnice?}
{\( \{ 10 \} \)}{\( \{ 1 \} \)}{\( \{ 100 \} \)}{\( \{ \ \} \)}{\( \left\{ \frac1{10} \right\} \)}{}}
\LANGsk{
\Question{
Je daná rovnica 
\[\log x+\log\sqrt x+\log\sqrt[4]x+\log\sqrt[8]x+\dots=2\]
s reálnou neznámou  \( x \). Aká je množina všetkých riešení tejto rovnice?}
{\( \{ 10 \} \)}{\( \{ 1 \} \)}{\( \{ 100 \} \)}{\( \{ \ \} \)}{\( \left\{ \frac1{10} \right\} \)}{}}
\LANGes{
\Question{
Dada la ecuación
\[\log x+\log\sqrt x+\log\sqrt[4]x+\log\sqrt[8]x+\dots=2\]
con una incógnita real \( x \). Calcula el conjunto de todas sus soluciones.}
{\( \{ 10 \} \)}{\( \{ 1 \} \)}{\( \{ 100 \} \)}{\( \{ \ \} \)}{\( \left\{ \frac1{10} \right\} \)}{}}
\End

\Data\ID{1003108808}
\Updated{2020-07-15 03:07:19}
\Level{B}
\SubArea{Infinite series}
\LANGen{
\Question{
We are given the equation
\[ \sum\limits_{n=1}^{\infty}2^{nx}=1 \]
with the unknown \( x \) being a real number. What is the set of all its solutions?}
{\( \{ -1 \} \)}{\( \{ 1 \} \)}{\( \{ 1;-1 \} \)}{\( \{\ \} \)}{\( \{ 0 \} \)}{}}
\LANGpl{
\Question{
Dane jest równanie
\[ \sum\limits_{n=1}^{\infty}2^{nx}=1 \]
z niewiadomą \( x \), która jest liczbą rzeczywistą. Wskaż zbiór jego wszystkich rozwiązań.}
{\( \{ -1 \} \)}{\( \{ 1 \} \)}{\( \{ 1;-1 \} \)}{\( \{\ \} \)}{\( \{ 0 \} \)}{}}
\LANGcs{
\Question{
Je dána rovnice
\[ \sum\limits_{n=1}^{\infty}2^{nx}=1 \]
s reálnou neznámou  \( x \). Jaká je množina všech řešení této rovnice?}
{\( \{ -1 \} \)}{\( \{ 1 \} \)}{\( \{ 1;-1 \} \)}{\( \{\ \} \)}{\( \{ 0 \} \)}{}}
\LANGsk{
\Question{
Je daná rovnica
\[ \sum\limits_{n=1}^{\infty}2^{nx}=1 \]
s reálnou neznámou  \( x \). Aká je množina všetkých riešení tejto rovnice?}
{\( \{ -1 \} \)}{\( \{ 1 \} \)}{\( \{ 1;-1 \} \)}{\( \{\ \} \)}{\( \{ 0 \} \)}{}}
\LANGes{
\Question{
Dada la ecuación
\[ \sum\limits_{n=1}^{\infty}2^{nx}=1 \]
con una incógnita real \( x \). Calcula el conjunto de todas sus soluciones.}
{\( \{ -1 \} \)}{\( \{ 1 \} \)}{\( \{ 1;-1 \} \)}{\( \{\ \} \)}{\( \{ 0 \} \)}{}}
\End

\Data\ID{1003108809}
\Updated{2020-07-15 03:07:19}
\Level{B}
\SubArea{Infinite series}
\LANGen{
\Question{
We are given the equation
\[ \sum\limits_{n=1}^{\infty} (\sin x)^{2n-2}=2\cdot\,\mathrm{tg}\,x \]
with the unknown \( x \) being a real number. What is the set of all its solutions?}
{\( \bigcup\limits_{k\in\mathbb{Z}}\left\{\frac14\pi+k\cdot\pi\right\} \)}{\( \bigcup\limits_{k\in\mathbb{Z}}\left\{\frac34\pi+k\cdot\pi\right\} \)}{\( \bigcup\limits_{k\in\mathbb{Z}}\left\{\frac34\pi+k\cdot\frac{\pi}2\right\} \)}{\( \bigcup\limits_{k\in\mathbb{Z}}\left\{\frac18\pi+k\cdot\frac{\pi}2\right\} \)}{\( \bigcup\limits_{k\in\mathbb{Z}}\left\{\frac14\pi+k\cdot\frac{\pi}2\right\} \)}{}}
\LANGpl{
\Question{
Dane jest równanie
\[ \sum\limits_{n=1}^{\infty} (\sin x)^{2n-2}=2\cdot\,\mathrm{tg}\,x \]
z niewiadomą \( x \), która jest liczbą rzeczywistą. Wskaż zbiór jego wszystkich rozwiązań.}
{\( \bigcup\limits_{k\in\mathbb{Z}}\left\{\frac14\pi+k\cdot\pi\right\} \)}{\( \bigcup\limits_{k\in\mathbb{Z}}\left\{\frac34\pi+k\cdot\pi\right\} \)}{\( \bigcup\limits_{k\in\mathbb{Z}}\left\{\frac34\pi+k\cdot\frac{\pi}2\right\} \)}{\( \bigcup\limits_{k\in\mathbb{Z}}\left\{\frac18\pi+k\cdot\frac{\pi}2\right\} \)}{\( \bigcup\limits_{k\in\mathbb{Z}}\left\{\frac14\pi+k\cdot\frac{\pi}2\right\} \)}{}}
\LANGcs{
\Question{
Je dána rovnice
\[ \sum\limits_{n=1}^{\infty} (\sin x)^{2n-2}=2\cdot\,\mathrm{tg}\,x \]
s reálnou neznámou  \( x \). Jaká je množina všech řešení této rovnice?}
{\( \bigcup\limits_{k\in\mathbb{Z}}\left\{\frac14\pi+k\cdot\pi\right\} \)}{\( \bigcup\limits_{k\in\mathbb{Z}}\left\{\frac34\pi+k\cdot\pi\right\} \)}{\( \bigcup\limits_{k\in\mathbb{Z}}\left\{\frac34\pi+k\cdot\frac{\pi}2\right\} \)}{\( \bigcup\limits_{k\in\mathbb{Z}}\left\{\frac18\pi+k\cdot\frac{\pi}2\right\} \)}{\( \bigcup\limits_{k\in\mathbb{Z}}\left\{\frac14\pi+k\cdot\frac{\pi}2\right\} \)}{}}
\LANGsk{
\Question{
Je daná rovnica
\[ \sum\limits_{n=1}^{\infty} (\sin x)^{2n-2}=2\cdot\,\mathrm{tg}\,x \]
s reálnou neznámou  \( x \). Aká je množina všetkých riešení tejto rovnice?}
{\( \bigcup\limits_{k\in\mathbb{Z}}\left\{\frac14\pi+k\cdot\pi\right\} \)}{\( \bigcup\limits_{k\in\mathbb{Z}}\left\{\frac34\pi+k\cdot\pi\right\} \)}{\( \bigcup\limits_{k\in\mathbb{Z}}\left\{\frac34\pi+k\cdot\frac{\pi}2\right\} \)}{\( \bigcup\limits_{k\in\mathbb{Z}}\left\{\frac18\pi+k\cdot\frac{\pi}2\right\} \)}{\( \bigcup\limits_{k\in\mathbb{Z}}\left\{\frac14\pi+k\cdot\frac{\pi}2\right\} \)}{}}
\LANGes{
\Question{
Dada la ecuación
\[ \sum\limits_{n=1}^{\infty} (\sin x)^{2n-2}=2\cdot\,\mathrm{tg}\,x \]
con una incógnita real \( x \). Calcula el conjunto de todas sus soluciones.}
{\( \bigcup\limits_{k\in\mathbb{Z}}\left\{\frac14\pi+k\cdot\pi\right\} \)}{\( \bigcup\limits_{k\in\mathbb{Z}}\left\{\frac34\pi+k\cdot\pi\right\} \)}{\( \bigcup\limits_{k\in\mathbb{Z}}\left\{\frac34\pi+k\cdot\frac{\pi}2\right\} \)}{\( \bigcup\limits_{k\in\mathbb{Z}}\left\{\frac18\pi+k\cdot\frac{\pi}2\right\} \)}{\( \bigcup\limits_{k\in\mathbb{Z}}\left\{\frac14\pi+k\cdot\frac{\pi}2\right\} \)}{}}
\End

\Data\ID{9000062903}
\Updated{2021-05-12 09:05:02}
\Level{B}
\SubArea{Infinite series}
\IMG{000629_03-figure0.pdf}
\LANGen{
\Question{
A semicircle is a half of the full circle. An infinite spiral is built from
semicircles with an increasing radius. The radius of the first semicircle is
\(3\, \mathrm{cm}\). The
radius of each semicircle in the spiral is bigger than the radius of the previous
semicircle by one third. Find the total length of the spiral.}
{\(\infty \)}{\(9\pi \)}{\(9\)}{\(3\pi \)}{}{}}
\LANGpl{
\Question{
Półkole to połowa okręgu. Nieskończona spirala składa się z półkoli o rosnącym promieniu. Promień pierwszego półkola jest równy 
\(3\, \mathrm{cm}\). 
Promień każdego półkola w spirali jest większy od promienia poprzedniego  o jedną trzecią.   Jaka będzie całkowita długość spirali?}
{\(\infty \)}{\(9\pi \)}{\(9\)}{\(3\pi \)}{}{}}
\LANGcs{
\Question{
„Nekonečná” spirála se skládá z půlkružnic. První půlkružnice má
poloměr 3 cm a každá další má poloměr o třetinu větší než
půlkružnice předcházející. Určete délku takto vzniklé spirály.}
{\(\infty \)}{\(9\pi \)}{\(9\)}{\(3\pi \)}{}{}}
\LANGsk{
\Question{
„Nekonečná” špirála sa skladá z polkružníc. Prvá polkružnica má
polomer 3 cm a každá ďalšia má polomer o tretinu väčší než
polkružnica predchádzajúca. Určte dĺžku takto vzniknutej špirály.}
{\(\infty \)}{\(9\pi \)}{\(9\)}{\(3\pi \)}{}{}}
\LANGes{
\Question{
Una espiral infinita está formada por semicircunferencias. El radio de la primera semicircunferencia mide
\(3\, \mathrm{cm}\). El radio de cada una de las semicircunferencias siguientes es un tercio del radio de la anterior. Calcula la longitud total de la espiral.}
{\(\infty \)}{\(9\pi \)}{\(9\)}{\(3\pi \)}{}{}}
\End

\Data\ID{9000062904}
\Updated{2021-05-12 09:05:17}
\Level{B}
\SubArea{Infinite series}
\IMG{000629_04-figure0.pdf}
\LANGen{
\Question{
A semicircle is a half of the full circle. An infinite spiral is built from
semicircles with a decreasing radius. The radius of the first semicircle is
\(3\, \mathrm{cm}\). The
radius of each semicircle in the spiral is smaller by one third of the radius of the
previous semicircle. Find the total length of the spiral.}
{\(9\pi \)}{\(9\)}{\(\frac{9}
{5}\pi \)}{\(\infty \)}{}{}}
\LANGpl{
\Question{
Nieskończona spirala zbudowana jest  z półkoli o malejącym promieniu. Promień pierwszego półkola jest równy 
\(3\, \mathrm{cm}\). 
Promień każdego półkola w spirali jest mniejszy o jedną trzecią od promienia poprzedniego półkola. Jaka będzie całkowita długość spirali?}
{\(9\pi \)}{\(9\)}{\(\frac{9}
{5}\pi \)}{\(\infty \)}{}{}}
\LANGcs{
\Question{
„Nekonečná” spirála se skládá z polokružnic. První polokružnice má
poloměr 3 cm a každá další má poloměr o třetinu menší než
polokružnice předcházející. Určete délku takto vzniklé spirály.}
{\(9\pi \)}{\(9\)}{\(\frac{9}
{5}\pi \)}{\(\infty \)}{}{}}
\LANGsk{
\Question{
„Nekonečná” špirála sa skladá z polkružníc. Prvá polkružnica má
polomer 3 cm a každá ďalšia má polomer o tretinu menší než
polkružnica predchádzajúca. Určte dĺžku takto vzniknutej špirály.}
{\(9\pi \)}{\(9\)}{\(\frac{9}
{5}\pi \)}{\(\infty \)}{}{}}
\LANGes{
\Question{
Una espiral infinita está formada por semicircunferencias. El radio de la primera semicircunferencia mide
\(3\, \mathrm{cm}\). El radio de cada una de las semicircunferencias siguientes es un tercio menor que el radio de la anterior. Calcula la longitud total de la espiral.}
{\(9\pi \)}{\(9\)}{\(\frac{9}
{5}\pi \)}{\(\infty \)}{}{}}
\End

\Data\ID{9000062905}
\Updated{2021-05-12 09:05:57}
\Level{B}
\SubArea{Infinite series}
\IMG{000629_05-figure0.pdf}
\LANGen{
\Question{
A semicircle is a half of the full circle. An infinite spiral is built from
semicircles with an increasing radius. The radius of the first semicircle is
\(2\, \mathrm{cm}\). The
radius of each semicircle in the spiral is a double of the radius of the previous
semicircle. Find the total length of the spiral.}
{\(\infty \)}{\(4\pi \)}{\(\frac{4}
{3}\pi \)}{\(- 4\pi \)}{}{}}
\LANGpl{
\Question{
Nieskończona spirala zbudowana jest z półkoli o rosnącym promieniu. Promień pierwszego półkola wynosi
\(2\, \mathrm{cm}\).
Promień każdego półkola w spirali jest dwukrotnością promienia poprzedniego półkola.  Jaka jest całkowita długość spirali?}
{\(\infty \)}{\(4\pi \)}{\(\frac{4}
{3}\pi \)}{\(- 4\pi \)}{}{}}
\LANGcs{
\Question{
„Nekonečná” spirála se skládá z půlkružnic. První půlkružnice má
poloměr 2 cm a každá další má poloměr dvakrát větší než
půlkružnice předcházející. Určete délku takto vzniklé spirály.}
{\(\infty \)}{\(4\pi \)}{\(\frac{4}
{3}\pi \)}{\(- 4\pi \)}{}{}}
\LANGsk{
\Question{
„Nekonečná” špirála sa skladá z polkružníc. Prvá polkružnica má
polomer 2 cm a každá ďalšia má polomer dvakrát väčší ako
polkružnica predchádzajúca. Určte dĺžku takto vzniknutej špirály.}
{\(\infty \)}{\(4\pi \)}{\(\frac{4}
{3}\pi \)}{\(- 4\pi \)}{}{}}
\LANGes{
\Question{
Una espiral infinita está formada por semicircunferencias. El radio de la primera semicircunferencia mide
\(2\, \mathrm{cm}\). El radio de cada una de las siguientes semicircunferencias en la espiral es el doble que el radio de la anterior. Calcula la longitud total de la espiral.}
{\(\infty \)}{\(4\pi \)}{\(\frac{4}
{3}\pi \)}{\(- 4\pi \)}{}{}}
\End

\Data\ID{9000062906}
\Updated{2021-05-12 09:05:51}
\Level{B}
\SubArea{Infinite series}
\IMG{000629_06-figure0.pdf}
\LANGen{
\Question{
A semicircle is a half of the full circle. An infinite spiral is built from
semicircles with a decreasing radius. The radius of the first semicircle is
\(2\, \mathrm{cm}\). The
radius of each semicircle in the spiral is one half of the radius of the previous
semicircle. Find the total length of the spiral.}
{\(4\pi \)}{\(\frac{4}
{3}\pi \)}{\(- 4\pi \)}{\(\infty \)}{}{}}
\LANGpl{
\Question{
Nieskończona spirala składa się z półkola o malejącym promieniu. Promień pierwszego półkola wynosi 
\(2\, \mathrm{cm}\).
Promień każdego półkola w spirali jest jedną drugą promienia poprzedniego półkola. Oblicz całkowitą długość spirali.}
{\(4\pi \)}{\(\frac{4}
{3}\pi \)}{\(- 4\pi \)}{\(\infty \)}{}{}}
\LANGcs{
\Question{
„Nekonečná” spirála se skládá z polokružnic. První polokružnice má
poloměr 2 cm a každá další má poloměr dvakrát menší než
polokružnice předcházející. Určete délku takto vzniklé spirály.}
{\(4\pi \)}{\(\frac{4}
{3}\pi \)}{\(- 4\pi \)}{\(\infty \)}{}{}}
\LANGsk{
\Question{
„Nekonečná” špirála sa skladá z polkružníc. Prvá polkružnica má
polomer 2 cm a každá ďalšia má polomer dvakrát menší ako
polkružnica predchádzajúca. Určte dĺžku takto vzniknutej špirály.}
{\(4\pi \)}{\(\frac{4}
{3}\pi \)}{\(- 4\pi \)}{\(\infty \)}{}{}}
\LANGes{
\Question{
Una espiral infinita está formada por semicircunferencias. El radio de la primera semicircunferencia mide
\(2\, \mathrm{cm}\). El radio de cada semicircunferencia siguiente en la espiral mide la mitad del radio de la anterior. Calcula la longitud total de la espiral.}
{\(4\pi \)}{\(\frac{4}
{3}\pi \)}{\(- 4\pi \)}{\(\infty \)}{}{}}
\End

\Data\ID{9000062907}
\Updated{2021-05-12 09:05:48}
\Level{B}
\SubArea{Infinite series}
\IMG{000629_07-figure0.pdf}
\LANGen{
\Question{
A quarter circle is an arc formed by one quarter of the full circle. An infinite spiral is
built from quarter circles with an increasing radius. The radius of the first quarter circle
is \(1\, \mathrm{cm}\).
The radius of each quarter circle in the spiral is bigger by one half than the radius of
the previous quarter circle. Find the total length of the spiral.}
{\(\infty \)}{\(4\pi \)}{\(\frac{2}
{5}\pi \)}{\(\frac{1}
{3}\pi \)}{}{}}
\LANGpl{
\Question{
Nieskończona spirala składa się z ćwierćkoli o rosnącym promieniu. Promień pierwszego ćwierćkoła wynosi 
\(1\, \mathrm{cm}\).
Promień każdego ćwierćkoła w spirali jest  większy o  jedną drugą od promienia poprzedniego ćwierćkoła. Oblicz całkowitą długość spirali.}
{\(\infty \)}{\(4\pi \)}{\(\frac{2}
{5}\pi \)}{\(\frac{1}
{3}\pi \)}{}{}}
\LANGcs{
\Question{
„Nekonečná” spirála se skládá ze čtvrtkružnic. První čtvrtkružnice má
poloměr 1 cm a každá další má poloměr o polovinu větší než
čtvrtkružnice předcházející. Určete délku takto vzniklé spirály.}
{\(\infty \)}{\(4\pi \)}{\(\frac{2}
{5}\pi \)}{\(\frac{1}
{3}\pi \)}{}{}}
\LANGsk{
\Question{
„Nekonečná” špirála sa skladá zo štvrťkružníc. Prvá štvrťkružnica má
polomer 1 cm a každá ďalšia má polomer o polovicu väčší než
štvrťkružnica predchádzajúca. Určte dĺžku takto vzniknutej špirály.}
{\(\infty \)}{\(4\pi \)}{\(\frac{2}
{5}\pi \)}{\(\frac{1}
{3}\pi \)}{}{}}
\LANGes{
\Question{
Una espiral infinita está formada por cuartos de circunferencias. El radio del primer cuarto de circunferencia mide
 \(1\, \mathrm{cm}\).
El radio de cada cuarto de circunferencia siguiente en la espiral es más grande por la mitad que el radio del anterior. Calcula la longitud total de la espiral.}
{\(\infty \)}{\(4\pi \)}{\(\frac{2}
{5}\pi \)}{\(\frac{1}
{3}\pi \)}{}{}}
\End

\Data\ID{9000062908}
\Updated{2021-05-12 10:05:19}
\Level{B}
\SubArea{Infinite series}
\IMG{000629_08-figure0.pdf}
\LANGen{
\Question{
A quarter circle is an arc formed by one quarter of the full circle. An infinite spiral is
built from quarter circles with an increasing radius. The radius of the first quarter circle
is \(4\, \mathrm{cm}\).
The radius of each quarter circle in the spiral is one half of the radius of the previous
quarter circle. Find the total length of the spiral.}
{\(4\pi \)}{\(8\)}{\(\frac{8}
{3}\)}{\(\infty \)}{}{}}
\LANGpl{
\Question{
Nieskończona spirala składa się z ćwierćkoli o rosnącym promieniu. Promień pierwszego ćwierćkoła wynosi 
\(4\, \mathrm{cm}\).
Promień każdego ćwierćkoła w spirali jest jedną drugą promienia poprzedniego ćwierćkoła. Oblicz całkowitą długość spirali.}
{\(4\pi \)}{\(8\)}{\(\frac{8}
{3}\)}{\(\infty \)}{}{}}
\LANGcs{
\Question{
„Nekonečná” spirála se skládá ze čtvrtkružnic. První čtvrtkružnice má
poloměr 4 cm a každá další má poloměr o polovinu menší než
čtvrtkružnice předcházející. Určete délku takto vzniklé spirály.}
{\(4\pi \)}{\(8\)}{\(\frac{8}
{3}\)}{\(\infty \)}{}{}}
\LANGsk{
\Question{
„Nekonečná” špirála sa skladá zo štvrťkružníc. Prvá štvrťkružnica má
polomer 4 cm a každá ďalšia má polomer o polovicu menší než
štvrťkružnica predchádzajúca. Určte dĺžku takto vzniknutej špirály.}
{\(4\pi \)}{\(8\)}{\(\frac{8}
{3}\)}{\(\infty \)}{}{}}
\LANGes{
\Question{
Una espiral infinita está formada por cuartos de circunferencias. El radio del primer cuarto de circunferencia mide \(4\, \mathrm{cm}\).
El radio de cada cuarto sigiente de circunferencia en la espiral mide la mitad del radio del cuarto anterior. Calcula la longitud total de la espiral.}
{\(4\pi \)}{\(8\)}{\(\frac{8}
{3}\)}{\(\infty \)}{}{}}
\End

\Data\ID{9000062909}
\Updated{2021-05-12 10:05:59}
\Level{B}
\SubArea{Infinite series}
\IMG{000629_09-figure0.pdf}
\LANGen{
\Question{
Consider the square of the side \(4\, \mathrm{cm}\).
The second square is inscribed into this first square by joining the centers of all sides.
In a similar way, the third square is inscribed into the second square by joining the
centers of the sides of the second square and this process continues up to infinity.
Find the sum of the perimeters of all squares.}
{\(32 + 16\sqrt{2}\)}{\(32 - 16\sqrt{2}\)}{\(32\)}{\(\infty \)}{}{}}
\LANGpl{
\Question{
Dany jest kwadrat o boku  \(4\, \mathrm{cm}\).
Drugi kwadrat został wpisany w pierwszy tak, by łączył środki jego boków. 
W podobny sposób wpisano kwadrat trzeci i kolejne aż do nieskończoności.
Oblicz sumę obwodów wszystkich kwadratów.}
{\(32 + 16\sqrt{2}\)}{\(32 - 16\sqrt{2}\)}{\(32\)}{\(\infty \)}{}{}}
\LANGcs{
\Question{
Je dán čtverec o straně 4 cm. Spojnice středů jeho stran tvoří opět
čtverec. Do tohoto čtverce je vepsán čtverec stejným způsobem atd.
Vypočítejte součet obvodů všech těchto čtverců.}
{\(32 + 16\sqrt{2}\)}{\(32 - 16\sqrt{2}\)}{\(32\)}{\(\infty \)}{}{}}
\LANGsk{
\Question{
Je daný štvorec so stranou dĺžky 4 cm. Spojnica stredov jeho strán tvorí opäť
štvorec. Do tohoto štvorca je vpísaný štvorec rovnakým spôsobom atď.
Vypočítajte súčet obvodov všetkých týchto štvorcov.}
{\(32 + 16\sqrt{2}\)}{\(32 - 16\sqrt{2}\)}{\(32\)}{\(\infty \)}{}{}}
\LANGes{
\Question{
Dado un cuadrado cuyo lado mide \(4\, \mathrm{cm}\).
Se inscribe un segundo cuadrado en el primer cuadrado uniendo los centros de todos los lados.
El proceso se repite de la misma manera. Halla la suma de los perímetros de todos los cuadrados.}
{\(32 + 16\sqrt{2}\)}{\(32 - 16\sqrt{2}\)}{\(32\)}{\(\infty \)}{}{}}
\End

\Data\ID{9000062910}
\Updated{2021-05-12 10:05:43}
\Level{B}
\SubArea{Infinite series}
\IMG{000629_10-figure0.pdf}
\LANGen{
\Question{
Consider the square of the side \(4\, \mathrm{cm}\).
The second square is inscribed into this first square by joining the centers of all sides.
In a similar way, the third square is inscribed into the second square by joining the
centers of the sides of the second square and this process continues up to infinity.
Find the sum of the squares of all squares.}
{\(32\)}{\(40\)}{\(\frac{32}
 {3} \)}{\(\infty \)}{}{}}
\LANGpl{
\Question{
Rozważ kwadrat o boku \(4\, \mathrm{cm}\).
Drugi kwadrat został wpisany w pierwszy tak, że łączy środki wszystkich jego boków.
W podobny sposób w drugi w
kwadrat wpisano trzeci i tak w nieskończoność.
Oblicz sumę kwadratów wszystkich kwadratów.}
{\(32\)}{\(40\)}{\(\frac{32}
 {3} \)}{\(\infty \)}{}{}}
\LANGcs{
\Question{
Je dán čtverec o straně 4 cm. Spojnice středů jeho stran tvoří opět
čtverec. Do tohoto čtverce je vepsán čtverec stejným způsobem atd.
Vypočítejte součet obsahů všech těchto čtverců.}
{\(32\)}{\(40\)}{\(\frac{32}
 {3} \)}{\(\infty \)}{}{}}
\LANGsk{
\Question{
Je daný štvorec so stranou dĺžky 4 cm. Spojnica stredov jeho strán tvorí opäť
štvorec. Do tohoto štvorca je vpísaný štvorec rovnakým spôsobom atď.
Vypočítajte súčet obsahov všetkých týchto štvorcov.}
{\(32\)}{\(40\)}{\(\frac{32}
 {3} \)}{\(\infty \)}{}{}}
\LANGes{
\Question{
Dado un cuadrado cuyo lado mide \(4\, \mathrm{cm}\).
Un segundo cuadrado se inscribe en el primer cuadrado uniendo los centros de todos los lados. El proceso se repite de la misma manera. Halla la suma de las áreas de todos los cuadrados.}
{\(32\)}{\(40\)}{\(\frac{32}
 {3} \)}{\(\infty \)}{}{}}
\End

\Data\ID{9000063407}
\Updated{2021-05-05 06:05:27}
\Level{B}
\SubArea{Infinite series}
\LANGen{
\Question{
In the following list find the value of the parameter
\(x\) which
ensure that the following geometric series is divergent. 
\[ 
 \sum _{n=1}^{\infty }(x + 4)^{2n}
\]}
{\(x = -5\)}{\(x = -\frac{9}
{2}\)}{\(x = -4\)}{\(x = -\frac{7}
{2}\)}{}{}}
\LANGpl{
\Question{
Z podanych odpowiedzi wybierz wartość parametru
\(x\),
która zapewni, że następujący szereg geometryczny jest rozbieżny.
\[ 
 \sum _{n=1}^{\infty }(x + 4)^{2n}
\]}
{\(x = -5\)}{\(x = -\frac{9}
{2}\)}{\(x = -4\)}{\(x = -\frac{7}
{2}\)}{}{}}
\LANGcs{
\Question{
Je dána nekonečná geometrická řada
\(\sum _{n=1}^{\infty }(x + 4)^{2n}\). Pro
které \(x\in \mathbb{R}\) je
tato řada divergentní?}
{\(x = -5\)}{\(x = -\frac{9}
{2}\)}{\(x = -4\)}{\(x = -\frac{7}
{2}\)}{}{}}
\LANGsk{
\Question{
Je daný nekonečný geometrický rad
\(\sum _{n=1}^{\infty }(x + 4)^{2n}\). Pre
ktoré \(x\in \mathbb{R}\) je
tento rad divergentný?}
{\(x = -5\)}{\(x = -\frac{9}
{2}\)}{\(x = -4\)}{\(x = -\frac{7}
{2}\)}{}{}}
\LANGes{
\Question{
Dada la serie geométrica infinita:
\(\sum _{n=1}^{\infty }(x + 4)^{2n}\). ¿Para
qué valor de \(x\in \mathbb{R}\) la serie es divergente?}
{\(x = -5\)}{\(x = -\frac{9}
{2}\)}{\(x = -4\)}{\(x = -\frac{7}
{2}\)}{}{}}
\End

\Data\ID{9000063408}
\Updated{2021-05-05 06:05:33}
\Level{B}
\SubArea{Infinite series}
\LANGen{
\Question{
In the following list find the value of the parameter
\(x\) which
ensure that the following geometric series is divergent. 
\[ 
 \sum _{n=1}^{\infty }(5 - 3x)^{n}
\]}
{\(x = \frac{1} 
{2}\)}{\(x = \frac{13} 
 {9} \)}{\(x = \frac{11} 
 {6} \)}{\(x = \frac{5} 
{3}\)}{}{}}
\LANGpl{
\Question{
Z podanych odpowiedzi wybierz wartość parametru 
\(x\), dla którego podany szereg geometryczny jest rozbieżny.
\[ 
 \sum _{n=1}^{\infty }(5 - 3x)^{n}
\]}
{\(x = \frac{1} 
{2}\)}{\(x = \frac{13} 
 {9} \)}{\(x = \frac{11} 
 {6} \)}{\(x = \frac{5} 
{3}\)}{}{}}
\LANGcs{
\Question{
Je dána nekonečná geometrická řada
\(\sum _{n=1}^{\infty }(5 - 3x)^{n}\). Pro
které \(x\in \mathbb{R}\) je
tato řada divergentní?}
{\(x = \frac{1} 
{2}\)}{\(x = \frac{13} 
 {9} \)}{\(x = \frac{11} 
 {6} \)}{\(x = \frac{5} 
{3}\)}{}{}}
\LANGsk{
\Question{
Je daný nekonečný geometrický rad
\(\sum _{n=1}^{\infty }(5 - 3x)^{n}\). Pre
ktoré \(x\in \mathbb{R}\) je
tento rad divergentný?}
{\(x = \frac{1} 
{2}\)}{\(x = \frac{13} 
 {9} \)}{\(x = \frac{11} 
 {6} \)}{\(x = \frac{5} 
{3}\)}{}{}}
\LANGes{
\Question{
Dada la serie geométrica infinita:
\(\sum _{n=1}^{\infty }(5 - 3x)^{n}\). ¿Para qué valor de \(x\in \mathbb{R}\) la serie es divergente?}
{\(x = \frac{1} 
{2}\)}{\(x = \frac{13} 
 {9} \)}{\(x = \frac{11} 
 {6} \)}{\(x = \frac{5} 
{3}\)}{}{}}
\End

\Data\ID{9000063409}
\Updated{2020-07-15 03:07:37}
\Level{B}
\SubArea{Infinite series}
\LANGen{
\Question{
Solve the following equation.
\[ 
 1 + 2x + 4x^{2} + 8x^{3}+\cdots = 3
\]}
{\(x = \frac{1} 
{3}\)}{\(x = \frac{1} 
{5}\)}{\(x = \frac{1} 
{2}\)}{\(x = \frac{3} 
{4}\)}{}{}}
\LANGpl{
\Question{
Rozwiąż następujące równanie.
\[ 
 1 + 2x + 4x^{2} + 8x^{3}+\cdots = 3
\]}
{\(x = \frac{1} 
{3}\)}{\(x = \frac{1} 
{5}\)}{\(x = \frac{1} 
{2}\)}{\(x = \frac{3} 
{4}\)}{}{}}
\LANGcs{
\Question{
Řešením rovnice \(1 + 2x + 4x^{2} + 8x^{3}+\cdots = 3\) 
je číslo:}
{\(x = \frac{1} 
{3}\)}{\(x = \frac{1} 
{5}\)}{\(x = \frac{1} 
{2}\)}{\(x = \frac{3} 
{4}\)}{}{}}
\LANGsk{
\Question{
Riešením rovnice \(1 + 2x + 4x^{2} + 8x^{3}+\cdots = 3\) 
je číslo:}
{\(x = \frac{1} 
{3}\)}{\(x = \frac{1} 
{5}\)}{\(x = \frac{1} 
{2}\)}{\(x = \frac{3} 
{4}\)}{}{}}
\LANGes{
\Question{
Resuelve la siguiente ecuación:
\[ 
 1 + 2x + 4x^{2} + 8x^{3}+\cdots = 3
\]}
{\(x = \frac{1} 
{3}\)}{\(x = \frac{1} 
{5}\)}{\(x = \frac{1} 
{2}\)}{\(x = \frac{3} 
{4}\)}{}{}}
\End

\Data\ID{9000063410}
\Updated{2020-07-15 03:07:37}
\Level{B}
\SubArea{Infinite series}
\LANGen{
\Question{
Solve the following equation. 
\[ 
 x + \frac{x} 
{3} + \frac{x} 
{9} + \frac{x} 
{27}+\cdots = 18
\]}
{\(x = 12\)}{\(x = 6\)}{\(x = 18\)}{\(x = 24\)}{}{}}
\LANGpl{
\Question{
Rozwiąż równanie.
\[ 
 x + \frac{x} 
{3} + \frac{x} 
{9} + \frac{x} 
{27}+\cdots = 18
\]}
{\(x = 12\)}{\(x = 6\)}{\(x = 18\)}{\(x = 24\)}{}{}}
\LANGcs{
\Question{
Řešením rovnice \(x + \frac{x} 
{3} + \frac{x} 
{9} + \frac{x} 
{27}+\cdots = 18\) 
je číslo:}
{\(x = 12\)}{\(x = 6\)}{\(x = 18\)}{\(x = 24\)}{}{}}
\LANGsk{
\Question{
Riešením rovnice \(x + \frac{x} 
{3} + \frac{x} 
{9} + \frac{x} 
{27}+\cdots = 18\) 
je číslo:}
{\(x = 12\)}{\(x = 6\)}{\(x = 18\)}{\(x = 24\)}{}{}}
\LANGes{
\Question{
Resuelve la siguiente ecuación:
\[ 
 x + \frac{x} 
{3} + \frac{x} 
{9} + \frac{x} 
{27}+\cdots = 18
\]}
{\(x = 12\)}{\(x = 6\)}{\(x = 18\)}{\(x = 24\)}{}{}}
\End

\Data\ID{9000073401}
\Updated{2020-07-15 03:07:37}
\Level{B}
\SubArea{Infinite series}
\LANGen{
\Question{
In the following list identify the expression which equals
\(3.3\overline{12}\).}
{\(3.3 +\sum _{ n=1}^{\infty }12\cdot 10^{-2n-1}\)}{\(3 +\sum _{ n=1}^{\infty }312\cdot 10^{-2n-1}\)}{\(3 +\sum _{ n=1}^{\infty }312\cdot 10^{-3n}\)}{\(3.3 +\sum _{ n=1}^{\infty }12\cdot 10^{-3n}\)}{}{}}
\LANGpl{
\Question{
Z podanych odpowiedzi wybierz wyrażenie, które jest równe
\(3.3\overline{12}\).}
{\(3.3 +\sum _{ n=1}^{\infty }12\cdot 10^{-2n-1}\)}{\(3 +\sum _{ n=1}^{\infty }312\cdot 10^{-2n-1}\)}{\(3 +\sum _{ n=1}^{\infty }312\cdot 10^{-3n}\)}{\(3.3 +\sum _{ n=1}^{\infty }12\cdot 10^{-3n}\)}{}{}}
\LANGcs{
\Question{
Určete, který z následujících výrazů se rovná číslu
\(3{,}3\overline{12}\).}
{\(3{,}3 +\sum _{ n=1}^{\infty }12\cdot 10^{-2n-1}\)}{\(3 +\sum _{ n=1}^{\infty }312\cdot 10^{-2n-1}\)}{\(3 +\sum _{ n=1}^{\infty }312\cdot 10^{-3n}\)}{\(3{,}3 +\sum _{ n=1}^{\infty }12\cdot 10^{-3n}\)}{}{}}
\LANGsk{
\Question{
Určte, ktorý z nasledujúcich výrazov sa rovná číslu
\(3{,}3\overline{12}\).}
{\(3{,}3 +\sum _{ n=1}^{\infty }12\cdot 10^{-2n-1}\)}{\(3 +\sum _{ n=1}^{\infty }312\cdot 10^{-2n-1}\)}{\(3 +\sum _{ n=1}^{\infty }312\cdot 10^{-3n}\)}{\(3{,}3 +\sum _{ n=1}^{\infty }12\cdot 10^{-3n}\)}{}{}}
\LANGes{
\Question{
Identifica la expresión que equivale a
\(3.3\overline{12}\).}
{\(3.3 +\sum _{ n=1}^{\infty }12\cdot 10^{-2n-1}\)}{\(3 +\sum _{ n=1}^{\infty }312\cdot 10^{-2n-1}\)}{\(3 +\sum _{ n=1}^{\infty }312\cdot 10^{-3n}\)}{\(3.3 +\sum _{ n=1}^{\infty }12\cdot 10^{-3n}\)}{}{}}
\End

\Data\ID{9000073402}
\Updated{2020-07-15 03:07:37}
\Level{B}
\SubArea{Infinite series}
\LANGen{
\Question{
In the following list identify the expression which equals
\(- 1.0\overline{345}\).}
{\(- 1 -\sum _{n=1}^{\infty }345\cdot 10^{-3n-1}\)}{\(- 1 -\sum _{n=1}^{\infty }345\cdot 10^{-3n}\)}{\(-\sum _{n=1}^{\infty }(10 + 345\cdot 10^{-3n-1})\)}{\(1 -\sum _{n=1}^{\infty }345\cdot 10^{-3n}\)}{}{}}
\LANGpl{
\Question{
Z podanych odpowiedzi wybierz wyrażenie, które jest równe
\(- 1.0\overline{345}\).}
{\(- 1 -\sum _{n=1}^{\infty }345\cdot 10^{-3n-1}\)}{\(- 1 -\sum _{n=1}^{\infty }345\cdot 10^{-3n}\)}{\(-\sum _{n=1}^{\infty }(10 + 345\cdot 10^{-3n-1})\)}{\(1 -\sum _{n=1}^{\infty }345\cdot 10^{-3n}\)}{}{}}
\LANGcs{
\Question{
Určete, který z následujících výrazů se rovná číslu
\(- 1{,}0\overline{345}\).}
{\(- 1 -\sum _{n=1}^{\infty }345\cdot 10^{-3n-1}\)}{\(- 1 -\sum _{n=1}^{\infty }345\cdot 10^{-3n}\)}{\(-\sum _{n=1}^{\infty }(10 + 345\cdot 10^{-3n-1})\)}{\(1 -\sum _{n=1}^{\infty }345\cdot 10^{-3n}\)}{}{}}
\LANGsk{
\Question{
Určte, ktorý z nasledujúcich výrazov sa rovná číslu
\(- 1{,}0\overline{345}\).}
{\(- 1 -\sum _{n=1}^{\infty }345\cdot 10^{-3n-1}\)}{\(- 1 -\sum _{n=1}^{\infty }345\cdot 10^{-3n}\)}{\(-\sum _{n=1}^{\infty }(10 + 345\cdot 10^{-3n-1})\)}{\(1 -\sum _{n=1}^{\infty }345\cdot 10^{-3n}\)}{}{}}
\LANGes{
\Question{
Identifica la expresión que equivale a 
\(- 1.0\overline{345}\).}
{\(- 1 -\sum _{n=1}^{\infty }345\cdot 10^{-3n-1}\)}{\(- 1 -\sum _{n=1}^{\infty }345\cdot 10^{-3n}\)}{\(-\sum _{n=1}^{\infty }(10 + 345\cdot 10^{-3n-1})\)}{\(1 -\sum _{n=1}^{\infty }345\cdot 10^{-3n}\)}{}{}}
\End

\Data\ID{9000073403}
\Updated{2020-07-15 03:07:37}
\Level{B}
\SubArea{Infinite series}
\LANGen{
\Question{
Find the sum of the following infinite series. 
\[ 
 -5\cdot 10^{-1} - 5\cdot 10^{-2} - 5\cdot 10^{-3} - 5\cdot 10^{-4}-\cdots 
\]}
{\(- 0.\overline{5}\)}{\(- 0.0\overline{5}\)}{\(0\)}{\(0.\overline{5}\)}{}{}}
\LANGpl{
\Question{
Oblicz sumę następujących szeregów nieskończonych.
\[ 
 -5\cdot 10^{-1} - 5\cdot 10^{-2} - 5\cdot 10^{-3} - 5\cdot 10^{-4}-\cdots 
\]}
{\(- 0.\overline{5}\)}{\(- 0.0\overline{5}\)}{\(0\)}{\(0.\overline{5}\)}{}{}}
\LANGcs{
\Question{
Určete, které z následujících desetinných čísel je rovno součtu nekonečné řady
\(- 5\cdot 10^{-1} - 5\cdot 10^{-2} - 5\cdot 10^{-3} - 5\cdot 10^{-4}-\cdots \).}
{\(- 0,\overline{5}\)}{\(- 0{,}0\overline{5}\)}{\(0\)}{\(0,\overline{5}\)}{}{}}
\LANGsk{
\Question{
Určte, ktoré z následujúcich desatinných čísel je rovné súčtu nekonečného radu
\(- 5\cdot 10^{-1} - 5\cdot 10^{-2} - 5\cdot 10^{-3} - 5\cdot 10^{-4}-\cdots \).}
{\(- 0,\overline{5}\)}{\(- 0{,}0\overline{5}\)}{\(0\)}{\(0,\overline{5}\)}{}{}}
\LANGes{
\Question{
Halla la suma de la siguiente serie infinita:
\[ 
 -5\cdot 10^{-1} - 5\cdot 10^{-2} - 5\cdot 10^{-3} - 5\cdot 10^{-4}-\cdots 
\]}
{\(- 0.\overline{5}\)}{\(- 0.0\overline{5}\)}{\(0\)}{\(0.\overline{5}\)}{}{}}
\End

\Data\ID{9000073407}
\Updated{2021-05-05 06:05:54}
\Level{B}
\SubArea{Infinite series}
\LANGen{
\Question{
Find all the values of \(x\) 
such that the following infinite series is convergent. 
\[ 
 1 + 3 - 2x + (3 - 2x)^{2} + (3 - 2x)^{3}+\cdots 
\]}
{\(x\in (1;2)\)}{\(x\in (-\infty ;-1)\)}{\(x\in (1;+\infty )\)}{\(x\in \mathbb{R}\)}{}{}}
\LANGpl{
\Question{
Dla jakich wartości parametru \(x\) 
podany szereg nieskończony jest zbieżny?
\[ 
 1 + 3 - 2x + (3 - 2x)^{2} + (3 - 2x)^{3}+\cdots 
\]}
{\(x\in (1;2)\)}{\(x\in (-\infty ;-1)\)}{\(x\in (1;+\infty )\)}{\(x\in \mathbb{R}\)}{}{}}
\LANGcs{
\Question{
Je dána nekonečná řada \(1 + 3 - 2x + (3 - 2x)^{2} + (3 - 2x)^{3}+\cdots \).
Určete, pro která \(x\) 
je řada konvergentní.}
{\(x\in (1;2)\)}{\(x\in (-\infty ;-1)\)}{\(x\in (1;+\infty )\)}{\(x\in \mathbb{R}\)}{}{}}
\LANGsk{
\Question{
Je daný nekonečný rad \(1 + 3 - 2x + (3 - 2x)^{2} + (3 - 2x)^{3}+\cdots \).
Určte, pre ktoré \(x\) 
je rad konvergentný.}
{\(x\in (1;2)\)}{\(x\in (-\infty ;-1)\)}{\(x\in (1;+\infty )\)}{\(x\in \mathbb{R}\)}{}{}}
\LANGes{
\Question{
Dada la serie infinita \(1 + 3 - 2x + (3 - 2x)^{2} + (3 - 2x)^{3}+\cdots \).
Determina para qué valores de \(x\) 
la serie es convergente.}
{\(x\in (1;2)\)}{\(x\in (-\infty ;-1)\)}{\(x\in (1;+\infty )\)}{\(x\in \mathbb{R}\)}{}{}}
\End

\Data\ID{9000073408}
\Updated{2021-05-05 06:05:33}
\Level{B}
\SubArea{Infinite series}
\LANGen{
\Question{
Find the values of \(x\) 
which ensure that the following infinite series is convergent. 
\[ 
 \sum _{n=1}^{\infty }\log ^{n-1}x
\]}
{\(x\in \left ( \frac{1}
{10};10\right )\)}{\(x\in (1;+\infty )\)}{\(x\in (1;10)\)}{\(x\in \mathbb{R}^{+}\)}{}{}}
\LANGpl{
\Question{
Określ wartości  \(x\), dla których następujący szereg nieskończony jest zbieżny.
\[ 
 \sum _{n=1}^{\infty }\log ^{n-1}x
\]}
{\(x\in \left ( \frac{1}
{10};10\right )\)}{\(x\in (1;+\infty )\)}{\(x\in (1;10)\)}{\(x\in \mathbb{R}^{+}\)}{}{}}
\LANGcs{
\Question{
Je dána nekonečná řada \(\sum _{n=1}^{\infty }\log ^{n-1}x\).
Určete, pro která \(x\) 
je řada konvergentní.}
{\(x\in \left ( \frac{1}
{10};10\right )\)}{\(x\in (1;+\infty )\)}{\(x\in (1;10)\)}{\(x\in \mathbb{R}^{+}\)}{}{}}
\LANGsk{
\Question{
Je daný nekonečný rad \(\sum _{n=1}^{\infty }\log ^{n-1}x\).
Určte, pre ktoré \(x\) 
je rad konvergentný.}
{\(x\in \left ( \frac{1}
{10};10\right )\)}{\(x\in (1;+\infty )\)}{\(x\in (1;10)\)}{\(x\in \mathbb{R}^{+}\)}{}{}}
\LANGes{
\Question{
Dada la serie infinita \(\sum _{n=1}^{\infty }\log ^{n-1}x\).
Determina para qué valores de \(x\) 
 la serie es convergente.}
{\(x\in \left ( \frac{1}
{10};10\right )\)}{\(x\in (1;+\infty )\)}{\(x\in (1;10)\)}{\(x\in \mathbb{R}^{+}\)}{}{}}
\End
